%  Соболевские пространства — конспект лекций
%  Компиляция: pdflatex main.tex  (дважды, из-за ссылок и оглавления)

\documentclass[12pt,a4paper]{article}

%  Преамбула конспекта «Соболевские пространства»
%  Компилировать pdfLaTeX (кодировка utf8, шрифты T2A).

% Latin Modern не имеет T2A-начертаний (нет t2almr.fd), поэтому кириллица идёт
% через cm-super: семейство cmr в кодировке T2A.
\usepackage[T2A]{fontenc}
\usepackage[utf8]{inputenc}
\usepackage[english,russian]{babel}

\usepackage{amsmath,amssymb,amsthm,mathtools}
\usepackage{geometry}
\geometry{a4paper,margin=2.5cm}

\usepackage{tikz}
\usetikzlibrary{arrows.meta,patterns,decorations.pathmorphing,decorations.pathreplacing,shapes.geometric,calc}
\usepackage{pgfplots}
\usepgfplotslibrary{groupplots,fillbetween}
\pgfplotsset{compat=1.18}

\usepackage[most]{tcolorbox}
\tcbuselibrary{breakable}

\usepackage{enumitem}
\usepackage{booktabs}
\usepackage{hyperref}
\hypersetup{
  colorlinks=true,
  linkcolor=blue!55!black,
  citecolor=blue!55!black,
  urlcolor=blue!55!black,
  pdftitle={Соболевские пространства --- конспект}
}

%  Окружения для теорем
\theoremstyle{plain}
% Счётчики независимы: в лекции теоремы, леммы и утверждения нумеруются
% каждое своим рядом (Т1, Т2, ... / Л1, Л2, ...).
\newtheorem{theorem}{Теорема}
\newtheorem{proposition}{Утверждение}
\newtheorem{corollary}{Следствие}
\newtheorem*{corollary*}{Следствие}
\newtheorem{lemma}{Лемма}
\newtheorem{criterion}{Признак}
\newtheorem*{criterion*}{Признак}

\theoremstyle{definition}
\newtheorem{definition}{Определение}
\newtheorem{example}{Пример}

\theoremstyle{remark}
\newtheorem*{remark}{Замечание}
\newtheorem*{reminder}{Напоминание}

\renewcommand{\qedsymbol}{$\blacksquare$}

%  Пометки конспекта
%    \begin{idea} ... \end{idea}  — пояснения, которых не было на лекции
%    \begin{warn} ... \end{warn} — неразборчиво / описка / требует сверки
\newtcolorbox{idea}[1][]{
  breakable,
  enhanced,
  colback=blue!4!white,
  colframe=blue!45!black,
  boxrule=0.4pt,
  left=6pt,right=6pt,top=5pt,bottom=5pt,
  title={Пояснение},
  fonttitle=\bfseries\small,
  coltitle=white,
  attach boxed title to top left={xshift=6pt,yshift=-2pt},
  boxed title style={colback=blue!45!black,boxrule=0pt,arc=1pt},
  #1
}

\newtcolorbox{warn}[1][]{
  breakable,
  enhanced,
  colback=orange!6!white,
  colframe=orange!65!black,
  boxrule=0.4pt,
  left=6pt,right=6pt,top=5pt,bottom=5pt,
  title={Сверить с записью},
  fonttitle=\bfseries\small,
  coltitle=white,
  attach boxed title to top left={xshift=6pt,yshift=-2pt},
  boxed title style={colback=orange!65!black,boxrule=0pt,arc=1pt},
  #1
}

% Отметка начала лекции
\newcommand{\lecture}[2]{%
  \bigskip
  \noindent\hfill\textit{\small #1\ifx\relax#2\relax\else, #2\fi}\\[-0.3em]
  \noindent\rule{\textwidth}{0.4pt}
  \medskip
}

%  Макросы
\newcommand{\R}{\mathbb{R}}
\newcommand{\N}{\mathbb{N}}
\newcommand{\Om}{\Omega}
\newcommand{\clOm}{\overline{\Omega}}
\newcommand{\dOm}{\partial\Omega}
\newcommand{\eps}{\varepsilon}
\newcommand{\dist}{\operatorname{dist}}
\newcommand{\supp}{\operatorname{supp}}
\newcommand{\esssup}{\operatorname*{ess\,sup}}
\newcommand{\Lp}[1]{L^{#1}}
\newcommand{\norm}[1]{\left\lVert #1 \right\rVert}
\newcommand{\abs}[1]{\left\lvert #1 \right\rvert}
\newcommand{\Lpn}[2]{\norm{#1}_{L^{#2}(\Omega)}}
\newcommand{\ball}[2]{B(#1,#2)}
\newcommand{\unif}{\rightrightarrows}
\newcommand{\Lloc}{L^1_{\mathrm{loc}}}
\newcommand{\Wp}[2]{W_{#1}^{#2}}
\newcommand{\Wpz}[2]{\mathring{W}_{#1}^{#2}}
\newcommand{\sgn}{\operatorname{sign}}
\newcommand{\dd}[2]{\frac{\partial #1}{\partial #2}}

% Маркер «++ в кружке» из тетради: доказательство будет дописано позже
% по записям прошлых семестров.
\newcommand{\plusplus}{%
  \tikz[baseline=(pp.base)]{%
    \node[draw, ellipse, inner xsep=3pt, inner ysep=1pt, line width=0.4pt] (pp)
      {\footnotesize $+\,+$};}%
}

%  Вспомогательные функции для рисунков
%    sstep(t)      — гладкая «ступенька» 6t^5-15t^4+10t^3, склеенная с 0 и 1
%    urho(x,rho)   — схематическое усреднение ступеньки u = 1_{[0,1)}
%    bump(x,rho)   — ядро усреднения omega_rho в размерности n = 1
\pgfmathdeclarefunction{sstep}{1}{%
  \pgfmathparse{(#1<0) ? 0 : ((#1>1) ? 1 : (6*(#1)^5-15*(#1)^4+10*(#1)^3))}%
}
\pgfmathdeclarefunction{urho}{2}{%
  \pgfmathparse{sstep((#1+#2)/(2*#2)) - sstep((#1-1+#2)/(2*#2))}%
}
% У края носителя 1-(x/rho)^2 -> 0, и деление 1/(...) переполняет pgfmath
% (предел ~16384) ещё до вызова exp. Поэтому знаменатель зажат снизу через max:
% показатель никогда не превосходит 15 по модулю, exp(-15) ~ 3e-7 — на графике ноль.
\pgfmathdeclarefunction{bump}{2}{%
  \pgfmathparse{(abs(#1)>=0.999*#2) ? 0 : (2.2523/#2)*exp(-1/max(1-((#1)/#2)^2,0.0667))}%
}

%  Быстрая сборка: \showfiguresfalse выключает все TikZ-рисунки.
%  Нужно, когда бесплатный Overleaf не успевает в лимит времени.
\newif\ifshowfigures
\showfigurestrue

\newcommand{\inputfig}[1]{%
  \ifshowfigures
    \input{#1}%
  \else
    \begin{center}\fbox{\small рисунок выключен: \texttt{\detokenize{#1}}}\end{center}%
  \fi
}


\title{\textbf{Соболевские пространства}\\[0.3em]\large Конспект лекций}
\author{}
\date{}

\begin{document}

\maketitle
\thispagestyle{empty}

\begin{idea}[title={Обозначения в конспекте}]
Блок \textbf{«Пояснение»} — то, чего не было на лекции: интуиция, обоснование
пропущенных шагов, связи с другими темами.

\medskip
Блок \textbf{«Сверить с записью»} — места, где запись на доске или в тетради
неразборчива, содержит описку или требует уточнения.

\medskip
Маркер \plusplus{} вместо доказательства — так помечено в тетради; эти доказательства
дописываются позже по записям прошлых семестров.
\end{idea}

\tableofcontents
\newpage

\lecture{Лекция 1}{}

\section*{Введение. Мотивирующая задача}
\addcontentsline{toc}{section}{Введение. Мотивирующая задача}

\subsection*{Можно ли восстановить функцию по её производным в точке?}

\textbf{Вопрос.} Пусть функция $f$ бесконечно дифференцируема в окрестности точки $x_0$ и
\[
  f(x_0) = 0, \qquad f^{(k)}(x_0) = 0 \quad \forall k \in \N .
\]
Чему равно $f(x_0 + \Delta x)$?

\textbf{Наивный ответ} — «нуль». Он получается из формулы Тейлора:
\[
  f(x_0 + \Delta x) = \sum_{k=0}^{\infty} \frac{f^{(k)}(x_0)}{k!}\,(\Delta x)^k = 0 .
\]

\textbf{Правильный ответ:} \emph{вообще говоря, ничего сказать нельзя}. Ряд Тейлора
действительно тождественно равен нулю, но гладкая функция \textbf{не обязана совпадать
со своим рядом Тейлора}. Совпадение гарантировано только для \emph{аналитических}
функций, а класс $C^\infty$ гораздо шире.

\begin{idea}
У аналитических функций (многочленов, $e^x$, $\sin x$ и т.\,п.) все производные в одной
точке определяют функцию в окрестности этой точки. Для функций из $C^\infty$ это неверно:
бесконечная гладкость — свойство гораздо более слабое, чем аналитичность. Именно эта
«гибкость» гладких функций и будет нам нужна.
\end{idea}

\subsection*{Контрпример: «функция-шапочка»}

\[
  f(x) =
  \begin{cases}
    e^{\frac{1}{x^2 - 1}}, & -1 < x < 1, \\[4pt]
    0, & |x| \ge 1 .
  \end{cases}
\]

\begin{warn}
На доске во второй строке было написано $x = \pm 1$. Из дальнейшего ясно, что имелось
в виду $|x| \ge 1$: функция продолжается нулём на всю прямую.
\end{warn}

Ту же функцию часто записывают в виде $f(x) = e^{-\frac{1}{1 - x^2}}$ при $|x| < 1$.

\textbf{Свойства:}
\begin{enumerate}
  \item \textbf{$f(x) > 0$ при $|x| < 1$}, например $f(0) = e^{-1}$.
  \item \textbf{$f$ непрерывна в точках $\pm 1$.} При $x \to \pm 1$ изнутри
        $\frac{1}{x^2-1} \to -\infty$, поэтому $f(x) \to 0$.
  \item \textbf{$f \in C^\infty(\R)$, причём $f^{(k)}(\pm 1) = 0$ для всех $k$.}
        По индукции производные внутри интервала имеют вид
        \[
          f^{(k)}(x) = \frac{P_k(x)}{(x^2-1)^{2k}}\; e^{\frac{1}{x^2-1}},
        \]
        где $P_k$ — многочлен. После замены $t = \frac{1}{1-x^2} \to +\infty$ получается
        выражение вида «многочлен от $t$, умноженный на $e^{-t}$», и оно стремится к $0$.
        Значит, все односторонние производные в точках $\pm 1$ нулевые и совпадают
        с производными нулевой функции снаружи.
\end{enumerate}

\textbf{Вывод.} В точке $x_0 = 1$ функция и все её производные равны нулю,
но $f(x_0 + \Delta x) \neq 0$ при $-2 < \Delta x < 0$.

\begin{idea}
\textbf{Зачем это в курсе.} Это пример \textbf{гладкой функции с компактным носителем}
(класс $C_c^\infty$, он же $C_0^\infty$ или $\mathcal{D}$). Аналитическая функция так себя
вести не может: если она равна нулю на интервале, то она равна нулю всюду. Функции из
$C_0^\infty$ — главный рабочий инструмент курса. Из «шапочки» строятся ядра усреднения
(глава~1), через $C_0^\infty$ определяются обобщённые производные.
\end{idea}


\newpage
% Нумерация: section = глава, subsection = параграф (§)
\renewcommand{\thesubsection}{\S\arabic{subsection}}

\section{Усреднение функций}

\subsection{Основные свойства}

\subsubsection*{Ядро усреднения}

\begin{definition}
Пусть $n \in \N$, $x \in \R^n$. Положим
\[
  \omega(x) =
  \begin{cases}
    C\, e^{\frac{1}{|x|^2 - 1}}, & |x| < 1, \\[4pt]
    0, & |x| \ge 1,
  \end{cases}
  \qquad
  C = \left( \int\limits_{\ball{0}{1}} e^{\frac{1}{|x|^2 - 1}}\, dx \right)^{-1} .
\]
\end{definition}

Здесь $|x| = \sqrt{x_1^2 + \dots + x_n^2}$, а $\ball{0}{1} = \{x : |x| < 1\}$ — открытый
единичный шар. Константа $C$ подобрана так, что
\[
  \int\limits_{\R^n} \omega(x)\, dx = \int\limits_{\ball{0}{1}} \omega(x)\, dx = 1 .
\]

\begin{idea}
\textbf{Почему $C$ определена корректно.} Подынтегральная функция непрерывна, положительна
в шаре и не превосходит $e^{-1}$, а шар имеет конечный объём. Поэтому интеграл конечен
и положителен.
\end{idea}

\textbf{Свойства $\omega$:}
\begin{enumerate}
  \item $\omega \ge 0$ на $\R^n$, причём $\omega > 0$ в $\ball{0}{1}$.
  \item $\supp \omega = \overline{\ball{0}{1}}$ — компакт.
  \item $\omega$ радиальна, то есть зависит только от $|x|$. В частности,
        $\omega(-x) = \omega(x)$.
  \item $\omega \in C^\infty(\R^n)$. Это видно из записи $\omega(x) = C\, g(|x|^2)$,
        где $g(t) = e^{\frac{1}{t-1}}$ при $t < 1$ и $g(t) = 0$ при $t \ge 1$. Функция $g$
        гладкая по тем же причинам, что и шапочка из введения, а
        $|x|^2 = x_1^2+\dots+x_n^2$ — многочлен. Запись через $|x|^2$ удобна тем, что сама
        норма $|x|$ в нуле не дифференцируема.
  \item $\displaystyle\int_{\R^n} \omega\,dx = 1$.
\end{enumerate}

\begin{remark}[из тетради]
Конкретный вид $\omega$ не важен: это просто бесконечно дифференцируемая финитная функция,
и вместо неё можно взять \textbf{любую другую такую же}.
\end{remark}

\begin{idea}
В доказательствах ниже используются только свойства 1, 2, 4 и 5: гладкость,
неотрицательность, носитель в единичном шаре и интеграл, равный единице. Любое ядро
с этими свойствами подходит.

\medskip
\textbf{Что это за объект.} $\omega$ — гладкая «плотность вероятности», сосредоточенная
в единичном шаре. Свёртка функции с таким ядром заменяет значение функции в точке на
взвешенное среднее по маленькой окрестности. Отсюда название — \textbf{усреднение}.
\end{idea}

\subsubsection*{Ядро с параметром $\rho$}

\begin{definition}
Пусть $\rho > 0$. Положим
\[
  \omega_\rho(x) = \rho^{-n}\, \omega\!\left(\frac{x}{\rho}\right),
  \qquad \supp \omega_\rho = \overline{\ball{0}{\rho}} .
\]
\end{definition}

\begin{warn}
На доске и в тетради написано $\supp\omega_\rho = \ball{0}{\rho}$. Носитель — замкнутое
множество, поэтому точнее писать $\overline{\ball{0}{\rho}}$. Для рассуждений это не важно.
\end{warn}

\textbf{Интеграл сохраняется:}
\[
  \int\limits_{\ball{0}{\rho}} \omega_\rho(x)\, dx
  = \int\limits_{\ball{0}{\rho}} \rho^{-n}\, \omega\!\left(\frac{x}{\rho}\right) dx
  \;\overset{\substack{y = x/\rho\\ dx = \rho^n dy}}{=}\;
  \int\limits_{\ball{0}{1}} \rho^{-n}\, \omega(y)\, \rho^n\, dy = 1 .
\]
Замена $y = x/\rho$ переводит шар $\ball{0}{\rho}$ в $\ball{0}{1}$, а её якобиан равен $\rho^n$.

Как и $\omega$, ядро $\omega_\rho$ гладкое, неотрицательное и радиальное.

\inputfig{figures/fig-kernel-rho}

\begin{idea}
\textbf{Картинка: уменьшение $\rho$.} Аргумент $x/\rho$ сжимает шапочку до шара радиуса
$\rho$, а множитель $\rho^{-n}$ поднимает её так, чтобы интеграл остался равным $1$.
Чем меньше $\rho$, тем шапочка уже и выше: $\omega_\rho(0) = C e^{-1}\rho^{-n} \to \infty$.
При $\rho \to 0$ вся «масса» собирается в нуле, то есть $\omega_\rho$ приближает
дельта-функцию.
\end{idea}

\subsubsection*{Усреднение функции}

Пусть $\Om \subset \R^n$ — открытое ограниченное множество (в тетради: «область»)
и $u \in \Lp{1}(\Om)$.

\begin{definition}
Функция
\[
  u_\rho(x) = \int\limits_{\Om} \omega_\rho(x - y)\, u(y)\, dy
\]
называется \textbf{усреднением функции $u$ с параметром $\rho > 0$}.
\end{definition}

\begin{idea}
При фиксированном $x$ функция $y \mapsto \omega_\rho(x-y)$ — шапочка с центром в точке $x$.
Поэтому $u_\rho(x)$ — взвешенное среднее значений $u$ в $\rho$-окрестности точки $x$.
По сути, $u_\rho = \omega_\rho * u$, где $u$ продолжена нулём вне $\Om$.
\end{idea}

\paragraph{Свойства усреднения.}

\textbf{1) $u_\rho : \R^n \to \R$, то есть $u_\rho$ определена во всех точках $\R^n$.}
Интеграл конечен, поскольку $\omega_\rho$ ограничена, а $u$ суммируема:
$|\omega_\rho(x-y)\,u(y)| \le \big(\sup \omega_\rho\big)\, |u(y)|$.
Подробнее об этом рассуждении — в отступлении ниже.

\textbf{2) Интегрирование фактически идёт по $\Om \cap \ball{x}{\rho}$.}
Поскольку $\supp_y\, \omega_\rho(x-y) = \overline{\ball{x}{\rho}}$,
\[
  u_\rho(x) = \int\limits_{\Om} \omega_\rho(x - y)\, u(y)\, dy
            = \int\limits_{\Om \cap \ball{x}{\rho}} \omega_\rho(x - y)\, u(y)\, dy ,
\]
и, в частности,
\[
  u_\rho(x) = 0, \quad \text{если } \dist(x, \Om) \ge \rho .
\]

\inputfig{figures/fig-support}

\begin{idea}
\textbf{Картинка из тетради.} Вокруг $\Om$ нарисована пунктиром $\rho$-окрестность. Если шар
$\ball{x}{\rho}$ целиком лежит внутри $\Om$, то это обычное усреднение. Если $x$ около
границы, шар частично выходит наружу, и интегрирование идёт только по части шара внутри
$\Om$; это всё равно что считать $u = 0$ вне $\Om$. Если $x$ дальше $\rho$ от $\Om$, шар
с $\Om$ не пересекается, и $u_\rho(x) = 0$. Итог: \textbf{носитель $u_\rho$ лежит
в замкнутой $\rho$-окрестности $\Om$}.
\end{idea}

\textbf{3) Гладкость: $u_\rho \in C^\infty(\R^n)$.} Интеграл «наследует» гладкость ядра
$\omega_\rho$: подынтегральная функция гладкая по $x$ и суммируемая по $y$. Используется
теорема о дифференцировании интеграла Лебега по параметру из курса анализа (как сказано
на лекции, она была в конце 5-го семестра).

\inputfig{figures/fig-smoothing}

\begin{idea}
\textbf{Подробнее.} Для любого мультииндекса $\alpha$
\[
  \big|D_x^\alpha\big[\omega_\rho(x-y)\,u(y)\big]\big|
  = \big|(D^\alpha\omega_\rho)(x-y)\big|\,|u(y)|
  \le \Big(\sup_{\R^n} |D^\alpha \omega_\rho|\Big)\, |u(y)| ,
\]
и правая часть — суммируемая мажоранта, не зависящая от $x$. Значит, дифференцировать
можно под знаком интеграла:
\[
  D^\alpha u_\rho(x) = \int\limits_\Om (D^\alpha \omega_\rho)(x-y)\, u(y)\, dy ,
\]
и эти производные непрерывны. Все производные «падают» на ядро, поэтому $u$ может быть
сколь угодно негладкой.

\medskip
\textbf{Картинка из тетради.} Изображён «рваный» график $u$, под ним маленькая шапочка,
скользящая вдоль оси. В каждой точке значение $u_\rho$ — среднее $u$ под шапочкой.
Поэтому $u_\rho$ сглаживает зубцы, и получается гладкая кривая.
\end{idea}

\subsubsection*{Отступление: мажорантный признак (вопрос с лекции)}

\begin{criterion}
Пусть $f, g$ измеримы на $E$, $|f| \le g$ на $E$ и $g$ суммируема. Тогда $f$ суммируема.
\end{criterion}

\begin{corollary*}
Если $u \in \Lp{1}(E)$ и $v \in \Lp{\infty}(E)$, то $uv \in \Lp{1}(E)$.
\end{corollary*}

Действительно, $uv$ измерима и
\[
  |uv| \le |u| \cdot \norm{v}_{\Lp{\infty}(E)} \quad \text{п.\,в. на } E .
\]
Здесь $|u|$ суммируема, а $\norm{v}_{\Lp{\infty}(E)} = \esssup_E |v|$ — константа.
Значит, $uv$ суммируема.

\begin{idea}
Именно этот признак работает в свойствах 1 и 3: $\omega_\rho$ и её производные ограничены
(лежат в $\Lp{\infty}$), а $u \in \Lp{1}$.
\end{idea}

\subsubsection*{Усреднение не увеличивает норму в $\Lp{p}$}

\begin{theorem}\label{thm:norm}
Пусть $1 \le p \le \infty$, $u \in \Lp{p}(\Om)$. Тогда
\[
  \norm{u_\rho}_{\Lp{p}(\Om)} \le \norm{u}_{\Lp{p}(\Om)} .
\]
\end{theorem}

\begin{idea}
Так как $\Om$ ограничено, $\Lp{p}(\Om) \subset \Lp{1}(\Om)$ (по неравенству Гёльдера
$\norm{u}_{\Lp{1}} \le |\Om|^{1-1/p}\norm{u}_{\Lp{p}}$). Поэтому усреднение $u_\rho$
определено.

\medskip
\textbf{Смысл.} Взвешенное среднее не может быть «больше» исходной функции. Иначе говоря,
оператор $u \mapsto u_\rho$ ограничен в $\Lp{p}$ и его норма не больше $1$.
\end{idea}

\textbf{Важное утверждение (используется во всех трёх случаях).} Для любого $x \in \R^n$
\begin{equation}\label{eq:star}
  \int\limits_{\Om} \omega_\rho(x - y)\, dy \le \int\limits_{\ball{x}{\rho}} \omega_\rho(x - y)\, dy = 1 .
\end{equation}
Неравенство верно, потому что $\omega_\rho \ge 0$ и
$\int_\Om = \int_{\Om\cap \ball{x}{\rho}} \le \int_{\ball{x}{\rho}}$. Последний интеграл
после замены $z = x - y$ превращается в $\int_{\ball{0}{\rho}}\omega_\rho(z)\,dz = 1$.

Тот же результат верен, если интегрировать по $x$, а не по $y$:
$\displaystyle\int_\Om \omega_\rho(x-y)\,dx \le 1$ при любом $y$.

\begin{proof}
\textbf{1$^\circ$. Случай $p = \infty$.}
\[
  |u_\rho(x)| \le \int\limits_\Om \omega_\rho(x - y)\, |u(y)|\, dy
  \le \norm{u}_{\Lp{\infty}(\Om)} \int\limits_\Om \omega_\rho(x - y)\, dy
  \overset{\eqref{eq:star}}{\le} \norm{u}_{\Lp{\infty}(\Om)} .
\]
На втором шаге вынесли $\esssup|u|$: $|u(y)| \le \norm{u}_{\Lp{\infty}}$ при почти всех $y$.
Следовательно, $\norm{u_\rho}_{\Lp{\infty}(\Om)} \le \norm{u}_{\Lp{\infty}(\Om)}$.

\begin{warn}
На доске замена записана как $z = y - x$, а в тетради — как $z = x - y$. Удобнее второй
вариант: при нём сразу $\omega_\rho(x-y) = \omega_\rho(z)$. При $z = y - x$ получается
$\omega_\rho(-z)$, и нужно дополнительно сослаться на чётность ядра. Также в тетради
в последнем интеграле написано $\omega(z)$ вместо $\omega_\rho(z)$ — это описка.
\end{warn}

\textbf{2$^\circ$. Случай $p = 1$.}
\begin{multline*}
  \int\limits_\Om |u_\rho(x)|\, dx
  \le \int\limits_\Om dx \int\limits_\Om \omega_\rho(x - y)\, |u(y)|\, dy
  \overset{\text{Фубини--Тонелли}}{=} \\
  = \int\limits_\Om |u(y)|\, dy
    \underbrace{\int\limits_\Om \omega_\rho(x - y)\, dx}_{\le 1 \text{ по } \eqref{eq:star}}
  \le \int\limits_\Om |u(y)|\, dy .
\end{multline*}
Значит, $\norm{u_\rho}_{\Lp{1}(\Om)} \le \norm{u}_{\Lp{1}(\Om)}$.

\begin{idea}
Менять порядок интегрирования можно по \textbf{теореме Тонелли}: подынтегральная функция
$\omega_\rho(x-y)|u(y)|$ неотрицательна, а для неотрицательных функций перестановка
допустима без предварительной проверки суммируемости.
\end{idea}

\textbf{3$^\circ$. Случай $1 < p < \infty$.} Пусть $q$ — сопряжённый показатель:
$\frac1p + \frac1q = 1$.

\begin{reminder}[неравенство Гёльдера]
Если $u \in \Lp{p}(\Om)$ и $v \in \Lp{q}(\Om)$, то $uv \in \Lp{1}(\Om)$ и
\[
  \int\limits_\Om |uv|\, dx \le \left(\int\limits_\Om |u|^p dx\right)^{1/p}
                                \left(\int\limits_\Om |v|^q dx\right)^{1/q} .
\]
\end{reminder}

Основной приём — разложить ядро на два множителя:
$\omega_\rho = \omega_\rho^{1/q}\cdot\omega_\rho^{1/p}$, поскольку $\frac1q + \frac1p = 1$.
Тогда по Гёльдеру
\begin{multline*}
  |u_\rho(x)| \le \int\limits_\Om \omega_\rho^{\frac1q}(x - y)\cdot \omega_\rho^{\frac1p}(x-y)\,|u(y)|\, dy \\
  \le \left( \int\limits_\Om \omega_\rho(x - y)\, dy \right)^{\frac1q}
      \left( \int\limits_\Om \omega_\rho(x - y)\, |u(y)|^p\, dy \right)^{\frac1p} .
\end{multline*}
Первая скобка не больше $1$ по \eqref{eq:star}. Возводя в степень $p$, получаем
\[
  |u_\rho(x)|^p \le \int\limits_\Om \omega_\rho(x - y)\, |u(y)|^p\, dy .
\]
Проинтегрируем по $x$ и снова применим Фубини--Тонелли и \eqref{eq:star}:
\begin{multline*}
  \norm{u_\rho}^p_{\Lp{p}(\Om)} = \int\limits_\Om |u_\rho(x)|^p dx
  \le \int\limits_\Om dx \int\limits_\Om \omega_\rho(x-y)\,|u(y)|^p\,dy \\
  = \int\limits_\Om |u(y)|^p\,dy \underbrace{\int\limits_\Om \omega_\rho(x-y)\,dx}_{\le 1}
  \le \norm{u}^p_{\Lp{p}(\Om)} . \qquad\qedhere
\end{multline*}
\end{proof}

\begin{idea}
\textbf{Идея случая 3$^\circ$.} Он сводится к случаю $p = 1$, применённому к функции
$|u|^p$. Разложение $\omega_\rho = \omega_\rho^{1/q}\omega_\rho^{1/p}$ нужно, чтобы
«забрать» часть веса в каждый сомножитель Гёльдера. Фактически это неравенство Йенсена
для выпуклой функции $t \mapsto |t|^p$ и вероятностной меры $\omega_\rho(x-y)\,dy$.
\end{idea}

\subsubsection*{Утверждение: равномерная сходимость на компактах}

\begin{proposition}\label{prop:unif}
Пусть $u \in C(\clOm)$, $K$ — компакт и $K \subset \Om$. Тогда
\[
  u_\rho \unif u \ \text{ на } K \quad \text{при } \rho \to 0 ,
\]
то есть $u_{\rho_k} \unif u$ на $K$ для любой последовательности $\rho_k > 0$, $\rho_k \to 0$.
\end{proposition}

\inputfig{figures/fig-compact}

\begin{proof}
\emph{Шаг 1: компакт отделён от границы.} Граница
$\dOm = \clOm \setminus \Om = \clOm \cap (\complement\Om)$ — всегда замкнутое множество.
Поскольку $\Om$ открыто, $\dOm \cap \Om = \varnothing$, а значит, $K \cap \dOm = \varnothing$.
Расстояние между непересекающимися компактом и замкнутым множеством положительно:
\[
  d = \dist(K, \dOm) > 0 .
\]

\emph{Шаг 2: при малых $\rho$ шары не вылезают из $\Om$.} Пусть $\rho < d$ и $x \in K$.
Тогда $\ball{x}{\rho} \subset \Om$, и поэтому
\[
  \int\limits_\Om \omega_\rho(x-y)\,dy
  = \int\limits_{\Om \cap \ball{x}{\rho}} \omega_\rho(x-y)\,dy
  = \int\limits_{\ball{x}{\rho}} \omega_\rho(x-y)\,dy = 1 .
\]

\begin{idea}
Почему $\ball{x}{\rho} \subset \Om$: если бы какая-то точка $z \in \ball{x}{\rho}$ не лежала
в $\Om$, то на отрезке $[x, z]$ нашлась бы точка границы $\dOm$ на расстоянии меньше
$\rho < d$ от $x \in K$. Это противоречит определению $d$. Обратите внимание: здесь
неравенство \eqref{eq:star} превращается в равенство. Около компакта усреднение берётся
по полному шару.
\end{idea}

\emph{Шаг 3: оценка разности.} Пусть $x \in K$, $\rho < d$. Представим $u(x)$ как
$u(x) \cdot 1$, заменив единицу интегралом из шага 2:
\[
  u(x) - u_\rho(x)
  = u(x)\int\limits_\Om \omega_\rho(x-y)\,dy - \int\limits_\Om \omega_\rho(x-y)\,u(y)\,dy
  = \int\limits_\Om \omega_\rho(x-y)\,\big(u(x) - u(y)\big)\,dy ,
\]
откуда
\[
  |u(x) - u_\rho(x)| \le \int\limits_{\ball{x}{\rho}} \omega_\rho(x-y)\,|u(x) - u(y)|\,dy .
\]

\emph{Шаг 4: равномерная непрерывность.} Возьмём $\eps > 0$. Найдётся $\delta > 0$ такое,
что $|u(x) - u(y)| < \eps$ при $|x - y| < \delta$. При $\rho < \min(d, \delta)$ все
$y \in \ball{x}{\rho}$ удовлетворяют $|x-y| < \delta$, поэтому
\[
  |u(x) - u_\rho(x)| \le \eps \int\limits_{\ball{x}{\rho}} \omega_\rho(x-y)\,dy = \eps
  \qquad \text{для всех } x \in K .
\]
Оценка не зависит от $x \in K$, значит, сходимость равномерная.
\end{proof}

\begin{warn}
В тетради написано «$u$ равномерно непрерывна на $K$». Этого недостаточно: точка $y$
пробегает шар $\ball{x}{\rho}$ и может не лежать в $K$. Правильно так: $\Om$ ограничено,
поэтому $\clOm$ — компакт, и $u \in C(\clOm)$ равномерно непрерывна на всём $\clOm$
(теорема Кантора).
\end{warn}

\begin{idea}
\textbf{Идея доказательства.} Вдали от границы усреднение — честное среднее по шару
с полным весом $1$. Поэтому $u - u_\rho$ — среднее отклонений $u(x)-u(y)$ по маленькому
шару, а по непрерывности эти отклонения малы.
\end{idea}

\begin{corollary*}
Если $u \in C(\clOm)$, то $u_\rho(x) \to u(x)$ при $\rho \to 0$ для каждого $x \in \Om$.
\end{corollary*}

\begin{proof}
Точка — это компакт: достаточно взять $K = \{x\}$.
\end{proof}

\subsubsection*{Сходимость усреднений в $\Lp{p}$}

\begin{theorem}\label{thm:conv}
Пусть $u \in \Lp{p}(\Om)$, $1 \le p < \infty$. Тогда
\[
  u_\rho \to u \ \text{ в } \Lp{p}(\Om) \quad \text{при } \rho \to 0 .
\]
\end{theorem}

\begin{idea}
\textbf{Почему $p < \infty$.} При $p = \infty$ теорема неверна. Все $u_\rho$ непрерывны,
а равномерный предел непрерывных функций непрерывен, поэтому разрывную функцию
(например, ступеньку из примера ниже) нельзя приблизить усреднениями в $\Lp{\infty}$.
\end{idea}

\begin{proof}
Сначала докажем для непрерывных функций, затем перенесём на все $\Lp{p}$ по плотности.

\textbf{1) Случай $u \in C(\clOm)$.} Возьмём компакт $K \subset \Om$ и разобьём
$\Om = K \cup (\Om \setminus K)$:
\[
  \norm{u - u_\rho}^p_{\Lp{p}(\Om)}
  = \int\limits_K |u - u_\rho|^p\,dx + \int\limits_{\Om\setminus K} |u - u_\rho|^p\,dx .
\]
Обозначим $M = \norm{u}_{C(\clOm)} = \norm{u}_{\Lp{\infty}(\Om)}$. По теореме~\ref{thm:norm}
(случай $p=\infty$) $\norm{u_\rho}_{\Lp{\infty}(\Om)} \le M$, поэтому $|u - u_\rho| \le 2M$
на $\Om$. Отсюда
\[
  \norm{u - u_\rho}^p_{\Lp{p}(\Om)}
  \le \sup_K |u - u_\rho|^p \cdot \mu(K) + (2M)^p \cdot \mu(\Om \setminus K) .
\]
Дальше по данному $\eps > 0$:
\begin{itemize}
  \item \textbf{сначала выбираем $K$} так, чтобы $(2M)^p\,\mu(\Om\setminus K) < \eps/2$.
        Это всегда возможно;
  \item \textbf{затем выбираем $\rho$}: так как $u_\rho \unif u$ на $K$
        (утверждение~\ref{prop:unif}), при всех достаточно малых $\rho$ выполнено
        $\sup_K |u - u_\rho|^p\,\mu(K) < \eps/2$.
\end{itemize}
Итого $\norm{u - u_\rho}^p_{\Lp{p}} < \eps$ для всех достаточно малых $\rho$.

\begin{idea}
\textbf{Почему $K$ можно выбрать.} Возьмём
$K_m = \{x \in \Om : \dist(x, \dOm) \ge \tfrac1m\}$. Это компакты (замкнуты и ограничены),
$K_m \subset \Om$, они возрастают и в объединении дают $\Om$. Поскольку $\mu(\Om) < \infty$,
по непрерывности меры $\mu(\Om \setminus K_m) \to 0$.

\medskip
\textbf{Порядок выбора важен:} $K$ выбирается раньше $\rho$ и от $\rho$ не зависит. Первое
слагаемое контролируется равномерной сходимостью на фиксированном компакте, второе —
малостью «каёмки» $\Om \setminus K$ около границы. Около границы сходимости может и не
быть, но там мало меры.
\end{idea}

\textbf{2) Общий случай $u \in \Lp{p}(\Om)$.} Пусть $\eps > 0$. Непрерывные функции плотны
в $\Lp{p}(\Om)$ при $1 \le p < \infty$, поэтому найдётся $g \in C(\clOm)$
с $\norm{u - g}_{\Lp{p}} < \eps/4$. Вставим $g$ и $g_\rho$:
\[
  \norm{u - u_\rho}_{\Lp{p}} = \norm{(u - g) + (g - g_\rho) + (g_\rho - u_\rho)}_{\Lp{p}}
  \le \norm{u - g}_{\Lp{p}} + \norm{g - g_\rho}_{\Lp{p}} + \norm{g_\rho - u_\rho}_{\Lp{p}} .
\]
Усреднение линейно по $u$, поэтому $g_\rho - u_\rho = (g - u)_\rho$. По теореме~\ref{thm:norm}
$\norm{(g-u)_\rho}_{\Lp{p}} \le \norm{g - u}_{\Lp{p}}$. Значит,
\[
  \norm{u - u_\rho}_{\Lp{p}}
  \le 2\underbrace{\norm{u - g}_{\Lp{p}}}_{< \eps/4}
   + \underbrace{\norm{g - g_\rho}_{\Lp{p}}}_{< \eps/2 \text{ при малых } \rho} < \eps .
\]
Второе слагаемое мало при малых $\rho$ по пункту 1, применённому к $g$.
\end{proof}

\begin{idea}
\textbf{Схема «$\eps/3$».} Это стандартный приём: (а) доказать сходимость на плотном классе
хороших функций; (б) иметь \textbf{равномерную по $\rho$} оценку нормы операторов, здесь —
теорему~\ref{thm:norm}. Тогда сходимость переносится на всё пространство. Теорема~\ref{thm:norm}
нужна именно здесь: без неё слагаемое $\norm{g_\rho - u_\rho}$ нечем было бы контролировать.

\medskip
Плотность: из курса теории меры известно, что в $\Lp{p}(\Om)$, $1 \le p < \infty$, плотны
даже непрерывные функции с компактным носителем в $\Om$. Продолженные нулём, они лежат
в $C(\clOm)$.
\end{idea}

\subsubsection*{Пример: усреднение ступеньки}

\begin{example}
Пусть $n = 1$, $\Om = (-1, 1)$,
\[
  u(x) = \begin{cases} 0, & x < 0, \\ 1, & x \ge 0, \end{cases}
  \qquad \rho \in \left\{1, \tfrac12, \tfrac14, \tfrac1{10}\right\}.
\]
Ядро: $\omega(x) = C e^{\frac{1}{x^2-1}}$ при $|x| < 1$ и $\omega(x) = 0$ при $|x| \ge 1$;
$\omega_\rho(x) = \frac1\rho\,\omega\!\left(\frac x\rho\right)$.
\end{example}

\begin{warn}
В тетради у $\omega$ пропущена нормирующая константа $C$. Для $n = 1$ получается
$C \approx 2{,}25$, так как $\int_{-1}^1 e^{\frac{1}{x^2-1}}dx \approx 0{,}444$. Например,
высота $\omega_{1/2}$ в нуле равна $2Ce^{-1} \approx 1{,}66$, а носитель — $[-\tfrac12, \tfrac12]$.
\end{warn}

\inputfig{figures/fig-step-setup}

Поскольку $u = 0$ на $(-1,0)$,
\[
  u_\rho(x) = \int\limits_{-1}^{1} \omega_\rho(x - y)\,u(y)\,dy = \int\limits_0^1 \omega_\rho(x - y)\,dy .
\]

\begin{idea}
\textbf{Как это посчитать.} После замены $z = x - y$ получаем
$u_\rho(x) = \int_{x-1}^{x} \omega_\rho(z)\,dz = \Phi_\rho(x) - \Phi_\rho(x-1)$, где
$\Phi_\rho(t) = \int_{-\infty}^t \omega_\rho$ — «гладкая ступенька»: $\Phi_\rho = 0$
при $t \le -\rho$, $\Phi_\rho = 1$ при $t \ge \rho$, а на $[-\rho, \rho]$ она гладко растёт.
Поэтому при $\rho \le \tfrac12$:

\medskip
\begin{center}
\begin{tabular}{ll}
\toprule
участок & значение $u_\rho(x)$ \\
\midrule
$x \le -\rho$ & $0$ \\
$-\rho \le x \le \rho$ & гладко растёт от $0$ до $1$ (в точке $x=0$ равно $\tfrac12$) \\
$\rho \le x \le 1-\rho$ & $1$ \\
$1-\rho \le x \le 1+\rho$ & гладко убывает от $1$ до $0$ (в точке $x = 1$ равно $\tfrac12$) \\
$x \ge 1+\rho$ & $0$ \\
\bottomrule
\end{tabular}
\end{center}
\end{idea}

\inputfig{figures/fig-step-mollified}

\paragraph{Графики из тетради.}
\begin{itemize}
  \item \textbf{Сверху:} $u_\rho$ около скачка в нуле. Вместо разрыва — гладкий подъём
        на отрезке $[-\rho, \rho]$, и чем меньше $\rho$, тем он круче.
  \item \textbf{Снизу:} $u_\rho$ на всей прямой.
        \begin{itemize}
          \item При $\rho = \tfrac12$ — «горб»: подъём на $[-\tfrac12, \tfrac12]$, значение
                $1$ ровно в точке $\tfrac12$, затем спуск до нуля к $x = \tfrac32$.
          \item При $\rho = \tfrac14$ — «трапеция»: подъём на $[-\tfrac14, \tfrac14]$,
                плато $u_\rho = 1$ на $[\tfrac14, \tfrac34]$, спуск на $[\tfrac34, \tfrac54]$.
        \end{itemize}
  \item При $\rho = 1$ функция $u_\rho$ вообще не достигает единицы: окно усреднения шире,
        чем участок, где $u = 1$.
\end{itemize}

\begin{idea}
\textbf{Что здесь видно.}
\begin{itemize}
  \item \textbf{Гладкость (свойство 3):} разрывная ступенька превратилась в бесконечно
        гладкую функцию.
  \item \textbf{Носитель (свойство 2):} $u_\rho$ «вылезает» за пределы $\Om$ на расстояние
        $\rho$ (правый спуск заканчивается в $1 + \rho$). Спуск около $x = 1$ возникает
        потому, что вне $\Om$ функция $u$ считается нулём.
  \item \textbf{Сходимость в $\Lp{p}$ (теорема~\ref{thm:conv}):} $u_\rho$ отличается от $u$
        только на отрезках длины порядка $\rho$ около $0$ и около $1$, и там
        $|u - u_\rho| \le 1$. Поэтому $\norm{u - u_\rho}^p_{\Lp{p}(\Om)} \lesssim \rho \to 0$.
  \item \textbf{Нет равномерной сходимости:} $u_\rho(0) = \tfrac12$ при всех малых $\rho$,
        а $u(0) = 1$. Утверждение о поточечной и равномерной сходимости требует
        непрерывности $u$, и ступенька показывает, что без неё оно неверно.
\end{itemize}
\end{idea}

\lecture{Лекция 2}{15.09.2026}

\subsubsection*{Следствия: плотность гладких функций в $\Lp{p}$}

\begin{corollary}\label{cor:dense-smooth}
$C^\infty(\clOm)$ плотно в $\Lp{p}(\Om)$, $1 \le p < \infty$. В качестве приближений можно
брать усреднения $u_\rho$.
\end{corollary}

\begin{proof}
По свойству 3 $u_\rho \in C^\infty(\R^n)$, поэтому сужение $u_\rho$ на $\Om$ лежит
в $C^\infty(\clOm)$: все производные непрерывно продолжаются до границы.
По теореме~\ref{thm:conv} $u_\rho \to u$ в $\Lp{p}(\Om)$.
\end{proof}

\begin{warn}
В тетради написано «$1 \le p = \infty$». Должно быть $1 \le p < \infty$: при $p = \infty$
утверждение неверно (см. пример со ступенькой).
\end{warn}

\begin{corollary}\label{cor:dense-c0}
$C_0^\infty(\Om)$ плотно в $\Lp{p}(\Om)$, $1 \le p < \infty$.
\end{corollary}

\begin{reminder}[что такое $C_0^\infty(\Om)$]
Если $u \in C_0^\infty(\Om)$, то, продолженная нулём, $u \in C^\infty(\R^n)$, и
\[
  \supp u \Subset \Om \iff \supp u \text{ — компакт и } \supp u \subset \Om .
\]
Знак $\Subset$ означает «строго внутри»: множество вместе с замыканием лежит в $\Om$.
\end{reminder}

\begin{idea}
\textbf{Разница со следствием~\ref{cor:dense-smooth}.} Функции из $C^\infty(\clOm)$ могут
быть ненулевыми вплоть до границы. Функции из $C_0^\infty(\Om)$ обязаны обращаться в ноль
в некоторой полосе около границы. Обычное усреднение $u_\rho$ этого не гарантирует: его
носитель, наоборот, вылезает за $\Om$. Поэтому функцию нужно сначала «обрезать» внутри $\Om$.
\end{idea}

\begin{proof}
Пусть $u \in \Lp{p}(\Om)$ и $\eps > 0$.

\emph{Шаг 1: обрезка.} Для компакта $K \subset \Om$ положим
\[
  u_K = u\,\chi_K = \begin{cases} u, & \text{на } K, \\ 0, & \text{вне } K. \end{cases}
\]
Тогда
\[
  \norm{u - u_K}_{\Lp{p}(\Om)}
  = \left(\int\limits_\Om |u - u_K|^p dx\right)^{1/p}
  = \left(\int\limits_{\Om\setminus K} |u - u_K|^p dx\right)^{1/p}
  = \left(\int\limits_{\Om\setminus K} |u|^p dx\right)^{1/p} .
\]
Поскольку $|u|^p \in \Lp{1}(\Om)$, по \textbf{абсолютной непрерывности интеграла Лебега}
(как функции множества) существует $\delta > 0$ такое, что
\[
  \int\limits_{E_0} |u|^p\,dx < \left(\frac\eps2\right)^p ,
  \quad \text{как только } \mu(E_0) < \delta .
\]
Выберем компакт $K \subset \Om$ с $\mu(\Om \setminus K) < \delta$; это возможно,
см. доказательство теоремы~\ref{thm:conv}. Тогда $\norm{u - u_K}_{\Lp{p}} < \eps/2$.

\emph{Шаг 2: усреднение обрезанной функции.} Рассмотрим
\[
  (u_K)_\rho(x) = \int\limits_\Om \omega_\rho(x-y)\,u_K(y)\,dy
                = \int\limits_{K \cap \ball{x}{\rho}} \omega_\rho(x-y)\,u_K(y)\,dy ,
\]
поскольку $u_K = 0$ вне $K$, а $\supp_y \omega_\rho(x-y) = \overline{\ball{x}{\rho}}$.
Если $\rho < \dist(K, \dOm)$, то
\[
  (u_K)_\rho(x) = 0 \quad \text{для всех } x \text{ с } \dist(x, K) \ge \rho ,
\]
то есть носитель $(u_K)_\rho$ лежит в замкнутой $\rho$-окрестности компакта $K$ — она сама
компакт и целиком содержится в $\Om$ (картинка в тетради: $K$ внутри $\Om$, вокруг $K$
пунктиром $\rho$-окрестность, ср. рис.~\ref{fig:compact}). Вместе с гладкостью (свойство 3)
это даёт
\[
  (u_K)_\rho \in C_0^\infty(\Om), \qquad \rho < \dist(K, \dOm) \quad\text{— финитная функция.}
\]

\emph{Шаг 3: оценка.} По теореме~\ref{thm:conv} $(u_K)_\rho \to u_K$ в $\Lp{p}(\Om)$,
поэтому найдётся $\rho_0 > 0$ такое, что $\norm{u_K - (u_K)_\rho}_{\Lp{p}} < \eps/2$
при $\rho < \rho_0$. Тогда
\[
  \norm{u - (u_K)_\rho}_{\Lp{p}}
  \le \norm{u - u_K}_{\Lp{p}} + \norm{u_K - (u_K)_\rho}_{\Lp{p}}
  < \frac\eps2 + \frac\eps2 = \eps
\]
для $\rho < \min\big(\rho_0,\ \dist(K, \dOm)\big)$.
\end{proof}

\begin{idea}
\textbf{Идея.} Сначала отрезаем от $u$ «каёмку» у границы — это стоит мало в $\Lp{p}$
из-за абсолютной непрерывности интеграла. Затем сглаживаем оставшуюся функцию усреднением
с радиусом меньше расстояния до границы. Сглаживание размывает носитель всего на $\rho$,
и он остаётся строго внутри $\Om$.
\end{idea}

\begin{remark}
Всё это — для $p < \infty$. В $\Lp{\infty}$ не плотны даже непрерывные функции.
\end{remark}

\subsubsection*{Локально суммируемые функции}

\begin{definition}
$u \in L^p_{\mathrm{loc}}(\Om)$ («$u$ локально принадлежит $\Lp{p}$»), если
\[
  u \in \Lp{p}(\Om') \quad \forall\, \Om' \Subset \Om .
\]
Случай $p = 1$ даёт класс $\Lloc(\Om)$ локально суммируемых функций — именно он нужен
дальше в курсе.
\end{definition}

\begin{example}
$\Om = \R$: функции, суммируемые на отрезках, лежат в $\Lloc{}(\R)$.
В частности, константа $\in \Lloc{}(\R)$, хотя она не суммируема на всей прямой.
\end{example}

\begin{example}
$\Om = (0,1)$: $\dfrac1x \in \Lloc{}\big((0,1)\big)$, хотя
$\int_0^1 \frac{dx}{x} = +\infty$ и $\tfrac1x \notin \Lp{1}((0,1))$.
\end{example}

\inputfig{figures/fig-loc-1x}

\begin{remark}[иерархия классов]
\[
  C_0^\infty(\Om) \text{ — самые хорошие функции (но самый узкий класс)},
\]
\[
  \Lloc{}(\Om) \text{ — самые плохие функции (но самый широкий класс)} .
\]
\end{remark}

\subsubsection*{Основная лемма вариационного исчисления (Дюбуа--Реймона)}

\begin{lemma}[Дюбуа--Реймон]\label{lem:dbr}
Пусть $\Om$ — открытое множество (не обязательно ограниченное),
\[
  f \in \Lloc{}(\Om), \qquad
  \int\limits_\Om f\,u\,dx = 0 \quad \forall u \in C_0^\infty(\Om) .
\]
Тогда $f = 0$ почти всюду на $\Om$.
\end{lemma}

\begin{idea}
Пометка с лекции: здесь основная борьба идёт за \emph{классы функций}. Пробные функции $u$
берутся из самого узкого класса $C_0^\infty$, а вывод делается про $f$ из самого широкого
класса $\Lloc{}$. Чем уже класс пробных функций, тем сильнее утверждение.
\end{idea}

\begin{proof}
\emph{Шаг 1: локализация.} Пусть $K \subset \Om$ — компакт. Найдётся открытое ограниченное
$\widetilde\Om$ с
\[
  K \subset \widetilde\Om \Subset \Om .
\]
Так как $f \in \Lloc{}(\Om)$, то $f \in \Lp{1}(\widetilde\Om)$.

\inputfig{figures/fig-nested}

\emph{Шаг 2: пробная функция.} Положим
\[
  g = \begin{cases} \operatorname{sign} f, & \text{на } K, \\ 0, & \text{вне } K, \end{cases}
  \qquad\text{то есть } g = \operatorname{sign} f \cdot \chi_K .
\]
Это произведение измеримых функций, поэтому $g$ измерима в $\Om$. Её усреднение
\[
  g_\rho(x) = \int\limits_{\widetilde\Om} \omega_\rho(x-y)\,g(y)\,dy
            = \int\limits_{K} \omega_\rho(x-y)\,g(y)\,dy ,
\]
и, как в следствии~\ref{cor:dense-c0},
\[
  g_\rho \in C_0^\infty(\widetilde\Om) \subset C_0^\infty(\Om),
  \qquad \text{если } \rho < \dist(K, \partial\widetilde\Om) .
\]

\emph{Шаг 3: сходимость почти всюду.} По теореме~\ref{thm:conv} $g_\rho \to g$
в $\Lp{1}(\widetilde\Om)$, значит, $g_\rho \to g$ по мере в $\widetilde\Om$. По теореме
Рисса существует последовательность $\rho_k \to 0$ с
\[
  g_{\rho_k} \to g \quad \text{п.\,в. в } \widetilde\Om
  \qquad\Longrightarrow\qquad
  f\,g_{\rho_k} \to f\,g \quad \text{п.\,в. в } \widetilde\Om .
\]

\emph{Шаг 4: мажоранта (нужная ветвь рассуждений).} По построению $g \in \Lp{\infty}(\widetilde\Om)$
и $|g| \le 1$ в $\widetilde\Om$. По теореме~\ref{thm:norm} (случай $p=\infty$)
\[
  \norm{g_\rho}_{\Lp{\infty}} \le \norm{g}_{\Lp{\infty}} \le 1
  \ \Longrightarrow\ |g_\rho| \le 1 \ \text{п.\,в. в } \widetilde\Om
  \ \Longrightarrow\ |g_\rho|\cdot|f| \le |f| \ \text{п.\,в. на } \widetilde\Om ,
\]
а $f \in \Lloc{}(\Om)$ даёт $|f| \in \Lp{1}(\widetilde\Om)$. Итого по теореме
Лебега о мажорированной сходимости
\[
  \int\limits_{\widetilde\Om} f\,g_{\rho_k}\,dx \longrightarrow \int\limits_{\widetilde\Om} f\,g\,dx .
\]

\emph{Шаг 5: предел равен нулю.} Но $g_{\rho_k} \in C_0^\infty(\Om)$, поэтому по условию леммы
\[
  \int\limits_{\widetilde\Om} f\,g_{\rho_k}\,dx = \int\limits_{\Om} f\,g_{\rho_k}\,dx = 0 .
\]
Следовательно,
\[
  0 = \int\limits_{\widetilde\Om} f\,g\,dx = \int\limits_K f\,\operatorname{sign} f\,dx
    = \int\limits_K |f|\,dx
  \qquad\Longrightarrow\qquad f = 0 \ \text{п.\,в. на } K .
\]

\emph{Шаг 6: от компакта ко всей области.} По теореме о представлении открытого множества
ячейками $\Om = \bigcup_{k=1}^{\infty} \overline{P_k}$, где $P_k$ — замкнутые ячейки, то есть
компакты. Тогда
\[
  \Om[f \neq 0] = \bigcup_{i \in \N} \overline{P_i}[f \neq 0]
  \ \Longrightarrow\
  \mu\big(\Om[f \neq 0]\big) \le \sum_{i=1}^{\infty} \mu\big(\overline{P_i}[f \neq 0]\big) = 0 ,
\]
поскольку каждое слагаемое нулевое по шагу 5. Значит, $f = 0$ п.\,в. на $\Om$.
\end{proof}

\begin{idea}
\textbf{Зачем $\operatorname{sign} f$.} Пробная функция подбирается так, чтобы
$f\,g = |f|$ на $K$: тогда равенство нулю интеграла сразу даёт $f = 0$. Но
$\operatorname{sign} f \cdot \chi_K$ разрывна и не годится как пробная функция, поэтому её
сначала сглаживают усреднением — отсюда переход к $g_\rho$ и предельный переход по Лебегу.
\end{idea}

\subsection{Критерий компактности в $\Lp{p}(\Om)$}

\begin{theorem}[критерий предкомпактности]\label{thm:compact}
Пусть $F \subset \Lp{p}(\Om)$, $1 \le p < \infty$, $\Om$ ограничено, и все функции из $F$
продолжены нулём на всё $\R^n$. Тогда $F$ предкомпактно $\iff$
\begin{enumerate}
  \item \textbf{равномерная ограниченность:} $\exists M \ge 0$ такое, что
        $\norm{u}_{\Lp{p}} \le M$ для всех $u \in F$;
  \item \textbf{равностепенная непрерывность сдвига:}
        \[
          \delta(\rho) \;:=\; \sup_{u \in F}\ \sup_{|z| < \rho}\
          \norm{u(x+z) - u(x)}_{\Lp{p}(\Om)} \xrightarrow[\rho \to 0]{} 0 ,
        \]
        где
        \[
          \norm{u(x+z) - u(x)}_{\Lp{p}(\Om)}
          = \left[\ \int\limits_\Om |u(x+z) - u(x)|^p\,dx \right]^{1/p} .
        \]
\end{enumerate}
\end{theorem}

\begin{warn}
Условие на $\Om$ на доске записано сокращённо (читается как «$\Om$ огр.»). В доказательстве
ограниченность $\Om$ используется дважды: в оценке $\mu(\Om)^{1/q}$ на шаге 1 и при переходе
от нормы в $C(\clOm)$ к норме в $\Lp{p}$ на шаге 3.
\end{warn}

\begin{idea}
Это критерий М.~Рисса--Колмогорова — аналог теоремы Арцела--Асколи для $\Lp{p}$. Условие 1 —
это «равномерная ограниченность», условие 2 — замена равностепенной непрерывности: функции
нельзя трясти по $x$ так, чтобы норма сдвига не стремилась к нулю равномерно по всему
семейству. Ключевая величина $\delta(\rho)$ зависит только от семейства $F$, не от отдельной
функции.
\end{idea}

\begin{proof}
Необходимость ($\Longrightarrow$) на лекции не разбиралась. Докажем достаточность
($\Longleftarrow$).

\paragraph{Шаг 1: при фиксированном $\rho$ семейство усреднений предкомпактно в $C(\clOm)$.}
Положим
\[
  F_\rho = \{\,u_\rho \mid u \in F\,\}, \qquad
  u_\rho(x) = \int\limits_\Om \omega_\rho(x-y)\,u(y)\,dy .
\]

\emph{Равномерная ограниченность $F_\rho$.} По неравенству Гёльдера
($\frac1p + \frac1q = 1$)
\begin{multline*}
  |u_\rho(x)|
  \le \left[\ \int\limits_{\Om \cap \ball{x}{\rho}} \omega_\rho^{\,q}(x-y)\,dy \right]^{1/q}
      \underbrace{\left[\ \int\limits_\Om |u(y)|^p\,dy \right]^{1/p}}_{\le\, M} \\
  \le M \left[\ \int\limits_{\ball{x}{\rho}} \omega_\rho^{\,q}(x-y)\,dy \right]^{1/q}
  = M\,\norm{\omega_\rho}_{\Lp{q}(\R^n)} .
\end{multline*}
Правая часть — некоторое число, зависящее только от $\rho$ (и не зависящее ни от $x$,
ни от $u$). Значит, $F_\rho$ равномерно ограничено.

\begin{idea}
Пометка с лекции: «но это при фиксированном $\rho$». Константа
$M\,\norm{\omega_\rho}_{\Lp{q}}$ растёт при $\rho \to 0$ (ядро становится выше), поэтому
равномерности по $\rho$ здесь нет и не требуется: на шаге 1 параметр $\rho$ фиксирован.
\end{idea}

\emph{Равностепенная непрерывность $F_\rho$.} Из
\[
  u_\rho(x+z) - u_\rho(x) = \int\limits_\Om \big(\omega_\rho(x+z-y) - \omega_\rho(x-y)\big)\,u(y)\,dy
\]
по неравенству Гёльдера
\[
  |u_\rho(x+z) - u_\rho(x)|
  \le \left[\ \int\limits_\Om \big|\omega_\rho(x+z-y) - \omega_\rho(x-y)\big|^q dy \right]^{1/q}
      \underbrace{\left[\ \int\limits_\Om |u(y)|^p dy \right]^{1/p}}_{\le\, M} .
\]
Ядро $\omega_\rho \in C^\infty$ и финитно, поэтому оно равномерно непрерывно на $\R^n$:
для любого $\eps > 0$ найдётся $\delta > 0$ такое, что
\[
  |z| < \delta \ \Longrightarrow\ \big|\omega_\rho(x+z-y) - \omega_\rho(x-y)\big| < \eps .
\]
Тогда
\[
  |u_\rho(x+z) - u_\rho(x)| \le \eps\,\big[\mu(\Om)\big]^{1/q}\,M
  \qquad \text{при } |z| < \delta ,
\]
причём оценка одна и та же для всех $x \in \Om$, всех $u_\rho \in F_\rho$ и всех $|z| < \delta$.
Значит, $F_\rho$ равностепенно непрерывно.

По теореме Арцела--Асколи семейство $F_\rho$ \textbf{предкомпактно в $C(\clOm)$ при каждом
фиксированном $\rho > 0$}.

\paragraph{Шаг 2: усреднения приближают $F$ равномерно по $u$.}
Вложим $\Om$ в большую ограниченную открытую область: $\Om \Subset \widetilde\Om$. Если
$\rho < \dist(\Om, \partial\widetilde\Om)$, то $\ball{x}{\rho} \subset \widetilde\Om$ для
всех $x \in \Om$.

\inputfig{figures/fig-omega-tilde}

Для $x \in \Om$
\[
  u(x) - u_\rho(x)
  = u(x) - \int\limits_\Om \omega_\rho(x-y)\,u(y)\,dy
  = u(x) - \int\limits_{\widetilde\Om} \omega_\rho(x-y)\,u(y)\,dy ,
\]
так как вне $\Om$ функция $u$ продолжена нулём. Далее, вес полный:
\[
  \int\limits_{\widetilde\Om} \omega_\rho(x-y)\,u(x)\,dy
  = \int\limits_{\widetilde\Om \cap \ball{x}{\rho}} \omega_\rho(x-y)\,u(x)\,dy
  = u(x)\underbrace{\int\limits_{\ball{x}{\rho}} \omega_\rho(x-y)\,dy}_{=\,1} = u(x) ,
\]
поэтому
\[
  u(x) - u_\rho(x) = \int\limits_{\widetilde\Om} \omega_\rho(x-y)\,\big(u(x) - u(y)\big)\,dy .
\]
По Гёльдеру с разложением $\omega_\rho = \omega_\rho^{1/q}\,\omega_\rho^{1/p}$:
\begin{multline*}
  |u(x) - u_\rho(x)|
  \le \underbrace{\left[\ \int\limits_{\widetilde\Om} \omega_\rho(x-y)\,dy \right]^{1/q}}_{\le\, 1}
      \left[\ \int\limits_{\widetilde\Om} \omega_\rho(x-y)\,|u(x) - u(y)|^p\,dy \right]^{1/p} \\
  \le \left[\ \int\limits_{\ball{0}{\rho}} \omega_\rho(z)\,|u(x) - u(x+z)|^p\,dz \right]^{1/p} ,
\end{multline*}
где использованы $\widetilde\Om \cap \ball{x}{\rho} = \ball{x}{\rho}$ и замена $z = y - x$.
Возведём в степень $p$ и проинтегрируем по $x$:
\begin{multline*}
  \int\limits_\Om |u(x) - u_\rho(x)|^p\,dx
  \le \int\limits_\Om dx \int\limits_{\ball{0}{\rho}} \omega_\rho(z)\,|u(x+z) - u(x)|^p\,dz \\
  \overset{\text{Фубини--Тонелли}}{=}
  \int\limits_{\ball{0}{\rho}} \omega_\rho(z)\,dz \int\limits_\Om |u(x+z) - u(x)|^p\,dx .
\end{multline*}
Внутренний интеграл не превосходит $\delta(\rho)^p$ по условию 2, а
$\int_{\ball{0}{\rho}} \omega_\rho = 1$, поэтому
\[
  \norm{u - u_\rho}_{\Lp{p}(\Om)} \le \delta(\rho) \xrightarrow[\rho \to 0]{} 0
  \qquad \forall u \in F .
\]
То есть $u_\rho \to u$ в $\Lp{p}$ \textbf{равномерно по $u \in F$}: оценка одна для всего
семейства.

\paragraph{Шаг 3: построение $\eps$-сети.}
Пусть $\eps > 0$. Выберем $\rho$ так, чтобы $\delta(\rho) < \eps/2$; тогда
\[
  \norm{u - u_\rho}_{\Lp{p}} \le \delta(\rho) < \frac\eps2 \qquad \forall u \in F .
\]
Положим
\[
  r = \frac{\eps}{2\,[\mu(\Om)]^{1/p}} .
\]
Семейство $F_\rho$ предкомпактно в $C(\clOm)$ по шагу 1, поэтому у него есть конечная
$r$-сеть $\{g_1, \dots, g_N\}$ в норме $C$. Для данного $u \in F$ возьмём
$u_\rho \in F_\rho$ и найдём номер $i$ с $\norm{u_\rho - g_i}_{C} < r$.
Переход от нормы $C$ к норме $\Lp{p}$ (здесь и нужна ограниченность $\Om$):
\[
  \norm{f}_{\Lp{p}} = \left(\int\limits_\Om |f|^p dx\right)^{1/p}
  \le [\mu(\Om)]^{1/p}\,\norm{f}_{C} ,
\]
поэтому $\norm{u_\rho - g_i}_{\Lp{p}} \le [\mu(\Om)]^{1/p}\,r = \eps/2$.
Итого
\[
  \norm{u - g_i}_{\Lp{p}}
  \le \underbrace{\norm{u - u_\rho}_{\Lp{p}}}_{<\,\eps/2}
    + \underbrace{\norm{u_\rho - g_i}_{\Lp{p}}}_{<\,\eps/2} < \eps ,
\]
то есть $\{g_1, \dots, g_N\}$ — конечная $\eps$-сеть семейства $F$ в $\Lp{p}(\Om)$.
Значит, $F$ предкомпактно.
\end{proof}

\begin{warn}
На доске шаг 3 сначала был записан через $\eps/3$-сеть ($\norm{u-u_\rho} < \eps/3$,
$\norm{u_\rho - g_i} < \eps/3$, итог $2\eps/3 < \eps$) — без перевода нормы из $C(\clOm)$
в $\Lp{p}(\Om)$. Затем он переписан аккуратно: сеть берётся с радиусом
$\eps/(2[\mu(\Om)]^{1/p})$ именно для того, чтобы после умножения на $[\mu(\Om)]^{1/p}$
получилось $\eps/2$. Выше приведён второй, корректный вариант.
\end{warn}

\begin{idea}
\textbf{Схема целиком.} Усреднение переводит $F$ в семейство $F_\rho$ гладких функций,
к которому применима Арцела--Асколи (шаг 1) — но только при фиксированном $\rho$. Условие 2
гарантирует, что $F_\rho$ приближает $F$ \emph{равномерно по всему семейству} (шаг 2).
Конечная сеть для $F_\rho$ превращается в конечную сеть для $F$ (шаг 3). Если бы в шаге 2
оценка зависела от $u$, сеть собрать бы не удалось — именно поэтому $\delta(\rho)$
определено как супремум по $u \in F$.
\end{idea}


\newpage
\section{Обобщённые функции и производные по Соболеву}

\subsection{Пространство основных функций $\mathcal{D}(\Om)$}

\subsubsection*{Мотивировка}

Пусть $u$ гладкая, $u \in C^1(a,b)$. Тогда для любой $\vartheta \in C_0^\infty\big((a,b)\big)$
интегрирование по частям даёт
\[
  \int\limits_a^b u'\,\vartheta\, dx = -\int\limits_a^b u\,\vartheta'\, dx
  \qquad \forall\, \vartheta \in C_0^\infty\big((a,b)\big) ,
\]
причём внеинтегральные члены пропадают: $\vartheta$ обращается в нуль около концов.
Правая часть имеет смысл и тогда, когда $u$ не дифференцируема, — её и берут
за определение производной.

\begin{idea}
Вся конструкция ниже — развитие этой одной формулы. Производную «перекладывают»
с плохой функции на хорошую пробную, расплачиваясь знаком $(-1)^{|\alpha|}$.
Поэтому сначала нужно аккуратно описать класс пробных функций и сходимость в нём.
\end{idea}

\subsubsection*{Основные функции и сходимость в $\mathcal{D}(\Om)$}

Пусть $\Om \subset \R^n$, $n \in \N$, не обязательно ограничено.

\begin{definition}
\[
  \mathcal{D}(\Om) = C_0^\infty(\Om)
\]
— множество \textbf{основных} (пробных) функций.
\end{definition}

\begin{definition}[сходимость в $\mathcal{D}$]
$\varphi_k \xrightarrow{\ \mathcal{D}\ } \varphi$, если
\begin{enumerate}
  \item $\exists\, K \subset \Om$ — компакт, такой что $\supp \varphi_k \subset K$
        для всех $k \in \N$ (все носители в одном компакте);
  \item $D^\alpha \varphi_k \unif D^\alpha \varphi$ равномерно, для любого мультииндекса
        $\alpha$.
\end{enumerate}
\end{definition}

\inputfig{figures/fig-test-conv}

$\mathcal{D}(\Om)$ — линейное пространство.

\begin{idea}
Требование общего компакта существенно: без него последовательность «шапочек»,
уползающих к границе $\Om$, формально сходилась бы к нулю равномерно вместе со всеми
производными, и непрерывность функционалов ничего бы не давала. Сходимость в
$\mathcal{D}$ намеренно сделана очень сильной, чтобы непрерывных функционалов было
как можно больше.
\end{idea}

\subsection{Распределения (обобщённые функции)}

\begin{definition}
$T$ — \textbf{распределение} (обобщённая функция) над $\Om$, если $T : \mathcal{D}(\Om) \to \R$
линеен и непрерывен:
\begin{enumerate}
  \item $T(a_1\varphi_1 + a_2\varphi_2) = a_1 T(\varphi_1) + a_2 T(\varphi_2)$
        для всех $\varphi_1, \varphi_2 \in \mathcal{D}(\Om)$, $a_1, a_2 \in \R$;
  \item $\varphi_k \xrightarrow{\ \mathcal{D}\ } \varphi
        \ \Longrightarrow\ T(\varphi_k) \to T(\varphi)$.
\end{enumerate}
Множество всех распределений над $\Om$ обозначается $\mathcal{D}^*(\Om)$.
\end{definition}

\subsubsection*{Регулярные распределения}

\begin{example}\label{ex:regular}
Пусть $f \in \Lloc(\Om)$. Положим
\[
  T_f(\varphi) = \int\limits_\Om f\,\varphi\, dx
  \qquad \forall\, \varphi \in \mathcal{D}(\Om) .
\]
Тогда $T_f \in \mathcal{D}^*(\Om)$.
\end{example}

\emph{Корректность.} Интеграл фактически берётся по $\supp \varphi$ — компакту внутри
$\Om$, а на компакте $f \in \Lp{1}(\supp\varphi)$ и $\varphi$ ограничена.

\emph{Линейность.}
\[
  T_f(a_1\varphi_1 + a_2\varphi_2) = \int\limits_\Om f\,(a_1\varphi_1 + a_2\varphi_2)\, dx
  = a_1 T_f(\varphi_1) + a_2 T_f(\varphi_2) .
\]

\emph{Непрерывность.} Пусть $\varphi_n \xrightarrow{\mathcal{D}} \varphi$ и $K$ — общий
компакт, содержащий все носители. Тогда
\[
  \left| \int\limits_K f\,\varphi_n\, dx - \int\limits_K f\,\varphi\, dx \right|
  \le \int\limits_K |f|\,|\varphi_n - \varphi|\, dx
  \le \sup_K |\varphi_n - \varphi| \int\limits_K |f|\, dx \xrightarrow[n \to \infty]{} 0 ,
\]
так как $\varphi_n \unif \varphi$, а $\int_K |f|\,dx < \infty$.

\begin{proposition}[единственность представления]\label{prop:reg-unique}
Если $f, g \in \Lloc(\Om)$ и $T_f(\varphi) = T_g(\varphi)$ для всех
$\varphi \in \mathcal{D}(\Om)$, то $f = g$ почти всюду на $\Om$.
\end{proposition}

\begin{proof}
$\int_\Om (f-g)\,\varphi\, dx = 0$ для всех $\varphi \in C_0^\infty(\Om)$, а это в точности
условие леммы Дюбуа--Реймона (лемма~\ref{lem:dbr}), откуда $f - g = 0$ п.\,в.
\end{proof}

\begin{definition}
Распределение $T \in \mathcal{D}^*(\Om)$ называется \textbf{регулярным}, если существует
$f \in \Lloc(\Om)$ с $T = T_f$, то есть
\[
  T(\varphi) = \int\limits_\Om f\,\varphi\, dx \qquad \forall \varphi \in \mathcal{D}(\Om) .
\]
Иначе распределение называется \textbf{сингулярным}.
\end{definition}

\begin{idea}
По утверждению~\ref{prop:reg-unique} функция $f$ в регулярном случае определена
однозначно (с точностью до множества меры нуль). Поэтому $\Lloc(\Om)$ вкладывается
в $\mathcal{D}^*(\Om)$, и распределения — честное расширение класса функций.
\end{idea}

\subsubsection*{$\delta$-функция Дирака}

\begin{example}[$\delta$-функция]\label{ex:delta}
Пусть $0 \in \Om$ и
\[
  \delta(\varphi) = \varphi(0) \qquad \forall \varphi \in \mathcal{D}(\Om) .
\]
Это распределение (линейность очевидна, непрерывность — из равномерной сходимости
$\varphi_k \unif \varphi$), причём \textbf{сингулярное}.
\end{example}

\begin{proof}[Доказательство сингулярности]
Пусть, от противного, $\delta = T_f$ для некоторой $f \in \Lloc(\Om)$, то есть
\[
  \int\limits_\Om f\,\varphi\, dx = \varphi(0) \qquad \forall \varphi \in \mathcal{D}(\Om) .
\]
Возьмём любую $\varphi \in C_0^\infty(\Om \setminus \{0\})$: для неё $\varphi(0) = 0$,
значит $\int_\Om f \varphi\, dx = 0$. По лемме Дюбуа--Реймона $f = 0$ п.\,в. на
$\Om \setminus \{0\}$, а точка имеет меру нуль, поэтому $f = 0$ п.\,в. на $\Om$. Но тогда
\[
  \varphi(0) = \int\limits_\Om f\,\varphi\, dx = 0 \qquad \forall \varphi \in \mathcal{D}(\Om) ,
\]
что неверно: достаточно взять «шапочку» с $\varphi(0) = e^{-1} \neq 0$. Противоречие.
\end{proof}

\begin{idea}
Наглядно $\delta$ — это предел ядер усреднения $\omega_\rho$ при $\rho \to 0$
(рис.~\ref{fig:kernel-rho}): все $\omega_\rho$ задают регулярные распределения
$T_{\omega_\rho}$, и $T_{\omega_\rho}(\varphi) \to \varphi(0)$, однако сам предел уже
не представляется никакой функцией.
\end{idea}

\subsection{Дифференцирование распределений}

\begin{definition}
Для $T \in \mathcal{D}^*(\Om)$ положим
\[
  D^\alpha T(\varphi) = (-1)^{|\alpha|}\, T(D^\alpha \varphi)
  \qquad \forall \varphi \in \mathcal{D}(\Om) .
\]
\end{definition}

\begin{proposition}[утв.\ 1]\label{prop:deriv-distr}
$D^\alpha T$ — тоже распределение.
\end{proposition}

\begin{proof}
\emph{Линейность} следует из линейности $T$ и оператора $D^\alpha$:
\[
  D^\alpha T(a_1\varphi_1 + a_2\varphi_2)
  = (-1)^{|\alpha|} T\big(a_1 D^\alpha\varphi_1 + a_2 D^\alpha\varphi_2\big)
  = a_1 D^\alpha T(\varphi_1) + a_2 D^\alpha T(\varphi_2) .
\]
\emph{Непрерывность}: если $\varphi_k \xrightarrow{\mathcal{D}} \varphi$, то и
$D^\alpha \varphi_k \xrightarrow{\mathcal{D}} D^\alpha \varphi$ (носители те же, а
равномерная сходимость требуется сразу для всех производных). Поэтому
\[
  D^\alpha T(\varphi_k) = (-1)^{|\alpha|} T(D^\alpha \varphi_k)
  \longrightarrow (-1)^{|\alpha|} T(D^\alpha \varphi) = D^\alpha T(\varphi) . \qedhere
\]
\end{proof}

\begin{corollary}
Любое распределение бесконечно дифференцируемо, причём
\[
  D^\alpha D^\beta T = D^{\alpha + \beta} T .
\]
\end{corollary}

\begin{proposition}[утв.\ 2: согласованность с классической производной]\label{prop:deriv-classic}
Если $f \in C^{|\alpha|}(\Om)$, то $D^\alpha T_f = T_{D^\alpha f}$.
\end{proposition}

\begin{proof}
\[
  D^\alpha T_f(\varphi) = (-1)^{|\alpha|} T_f(D^\alpha \varphi)
  = (-1)^{|\alpha|} \int\limits_\Om f\, D^\alpha \varphi\, dx
  = \int\limits_\Om D^\alpha f \cdot \varphi\, dx = T_{D^\alpha f}(\varphi) ,
\]
где предпоследнее равенство — интегрирование по частям $|\alpha|$ раз; внеинтегральные
члены равны нулю, поскольку $\varphi$ финитна.
\end{proof}

\begin{idea}
Это и есть оправдание всей конструкции: на гладких функциях новое дифференцирование
совпадает со старым, а область определения расширяется на весь $\mathcal{D}^*(\Om)$.
\end{idea}

\subsection{Обобщённая производная по Соболеву}

\begin{definition}\label{def:sobolev-deriv}
Пусть $f \in \Lloc(\Om)$. Функция $g \in \Lloc(\Om)$ называется \textbf{обобщённой
производной $f$ по Соболеву}, $g = D^\alpha f$, если
\[
  D^\alpha T_f = T_g ,
\]
то есть
\[
  \int\limits_\Om g\,\varphi\, dx = (-1)^{|\alpha|} \int\limits_\Om f\, D^\alpha \varphi\, dx
  \qquad \forall \varphi \in C_0^\infty(\Om) .
\]
\end{definition}

\begin{remark}
Производная по Соболеву существует \textbf{не всегда}: распределение $D^\alpha T_f$
определено для любой $f \in \Lloc(\Om)$, но может оказаться сингулярным — тогда
никакой функции $g$ ему не соответствует.
\end{remark}

\inputfig{figures/fig-heaviside}

\begin{example}[функция Хевисайда]\label{ex:heaviside}
Пусть $\Om = \R$ и
\[
  H(x) = \begin{cases} 0, & x < 0, \\ 1, & x \ge 0. \end{cases}
\]
Тогда
\begin{multline*}
  D T_H(\varphi) = -\,T_H(\varphi')
  = -\int\limits_{-\infty}^{+\infty} H(x)\,\varphi'(x)\, dx \\
  = -\int\limits_0^{+\infty} \varphi'(x)\, dx
  = -\big(0 - \varphi(0)\big) = \varphi(0) = \delta(\varphi) ,
\end{multline*}
то есть $H' = \delta$. Распределение $\delta$ сингулярно (пример~\ref{ex:delta}),
значит \textbf{у функции Хевисайда нет соболевской производной}.
\end{example}

\begin{example}\label{ex:x2sign}
Пусть
\[
  f(x) = x^2 \sgn x = \begin{cases} x^2, & x \ge 0, \\ -x^2, & x < 0. \end{cases}
\]
Тогда
\[
  f'(x) = \begin{cases} 2x, & x \ge 0, \\ -2x, & x < 0, \end{cases} = 2|x| ,
  \qquad
  f''(x) = \begin{cases} 2, & x \ge 0, \\ -2, & x < 0, \end{cases} = 2\sgn x ,
\]
и обе эти производные — обычные функции из $\Lloc(\R)$. А вот третья производная уже
сингулярна:
\[
  f'''(\varphi) = -\,T_{f''}(\varphi') = -\int\limits_{-\infty}^{+\infty} 2\sgn x \cdot \varphi'(x)\, dx
  = 2\int\limits_{-\infty}^{0}\varphi'\,dx - 2\int\limits_0^{+\infty}\varphi'\,dx
  = 4\varphi(0) ,
\]
то есть $f''' = 4\delta$.
\end{example}

\begin{idea}
Видна общая картина: каждое дифференцирование «съедает» одну степень гладкости, и как
только у функции появляется скачок, следующая производная выпадает из $\Lloc$ и становится
сингулярной. Функция из примера~\ref{ex:x2sign} лежит в $\Wp{p}{2}$, но не в $\Wp{p}{3}$.
\end{idea}

\subsubsection*{Свойства обобщённой производной}

\begin{proposition}[свойство 1: единственность]\label{prop:deriv-unique}
Обобщённая производная единственна с точностью до множества меры нуль: если $g_1$ и $g_2$
обе удовлетворяют определению~\ref{def:sobolev-deriv}, то $g_1 = g_2$ почти всюду в $\Om$.
\end{proposition}

\begin{proof}
Вычитая два равенства, получаем
\[
  \int\limits_\Om (g_1 - g_2)\,\varphi\, dx = 0 \qquad \forall \varphi \in C_0^\infty(\Om) ,
\]
откуда по лемме Дюбуа--Реймона (лемма~\ref{lem:dbr}) $g_1 - g_2 = 0$ п.\,в.
\end{proof}

\begin{proposition}[свойство 2: линейность]\label{prop:deriv-linear}
Пусть $f_1, f_2 \in \Lloc(\Om)$ и существуют $D^\alpha f_1, D^\alpha f_2 \in \Lloc(\Om)$.
Тогда для любых $a_1, a_2 \in \R$ существует
$D^\alpha(a_1 f_1 + a_2 f_2) \in \Lloc(\Om)$ и
\[
  D^\alpha(a_1 f_1 + a_2 f_2) = a_1 D^\alpha f_1 + a_2 D^\alpha f_2 .
\]
\end{proposition}

\begin{proof}
Складывая определяющие равенства с коэффициентами $a_1, a_2$, получаем для всех
$\varphi \in C_0^\infty(\Om)$
\[
  \int\limits_\Om \big(a_1 D^\alpha f_1 + a_2 D^\alpha f_2\big)\varphi\, dx
  = (-1)^{|\alpha|} \int\limits_\Om \big(a_1 f_1 + a_2 f_2\big) D^\alpha \varphi\, dx ,
\]
а это и есть определение~\ref{def:sobolev-deriv} для функции $a_1 f_1 + a_2 f_2$.
\end{proof}

\subsubsection*{Сходимость в $\Lloc(\Om)$}

\begin{definition}
$f_k \to f$ в $\Lloc(\Om)$, если
\[
  \int\limits_{\Om'} |f_k - f|\, dx \xrightarrow[k \to \infty]{} 0
  \qquad \forall\, \Om' \Subset \Om ,
\]
то есть сходимость в $\Lp{1}(\Om')$ на каждом компактно вложенном подмножестве.
\end{definition}

\begin{remark}[метрика]
Эта сходимость метризуема. Возьмём $\Om_k \Subset \Om$ с $\bigcup_k \Om_k = \Om$ и
положим
\[
  \rho_k(f_1, f_2) = \int\limits_{\Om_k} |f_1 - f_2|\, dx ,
  \qquad
  \rho(f_1, f_2) = \sum_{k=1}^{\infty} \frac{1}{2^k}\,
  \frac{\rho_k(f_1, f_2)}{1 + \rho_k(f_1, f_2)} .
\]
Ряд сходится, так как каждое слагаемое не превосходит $2^{-k}$, и сумма не больше $1$.
\end{remark}

\begin{idea}
Дробь $\frac{t}{1+t}$ обрезает каждое слагаемое до величины меньше единицы — без неё
ряд мог бы расходиться, ведь $\rho_k$ не ограничены. При этом $\frac{t}{1+t} \to 0$
тогда и только тогда, когда $t \to 0$, так что сходимость по метрике $\rho$ равносильна
сходимости всех $\rho_k$ к нулю, то есть сходимости на каждом $\Om' \Subset \Om$.
\end{idea}

\subsubsection*{Свойство 3: признак существования обобщённой производной}

\begin{theorem}[признак существования производной по Соболеву]\label{thm:crit-deriv}
Пусть
\[
  f_k \in \Lloc(\Om) \quad \forall k, \qquad
  \exists\, D^\alpha f_k \in \Lloc(\Om) \quad \forall k,
\]
\[
  f_k \to f \ \text{ в } \Lloc(\Om), \qquad
  D^\alpha f_k \to g \ \text{ в } \Lloc(\Om) .
\]
Тогда
\[
  \exists\, D^\alpha f \in \Lloc(\Om): \qquad D^\alpha f = g .
\]
\end{theorem}

\begin{proof}
Пусть $\varphi \in C_0^\infty(\Om)$. По определению обобщённой производной для каждого $k$
\[
  \int\limits_\Om f_k\, D^\alpha \varphi\, dx
  = (-1)^{|\alpha|} \int\limits_\Om D^\alpha f_k \cdot \varphi\, dx ,
\]
а так как $\varphi$ финитна, интегрирование идёт лишь по $\supp \varphi$:
\[
  \int\limits_{\supp \varphi} f_k\, D^\alpha \varphi\, dx
  = (-1)^{|\alpha|} \int\limits_{\supp \varphi} D^\alpha f_k \cdot \varphi\, dx .
\]

\emph{Предел в левой части.}
\begin{multline*}
  \left| \int\limits_{\supp \varphi} f_k\, D^\alpha \varphi\, dx
       - \int\limits_{\supp \varphi} f\, D^\alpha \varphi\, dx \right|
  \le \int\limits_{\supp \varphi} |f_k - f|\cdot |D^\alpha \varphi|\, dx \\
  \le \sup_{\R^n} |D^\alpha \varphi| \int\limits_{\supp \varphi} |f_k - f|\, dx
  \xrightarrow[k \to \infty]{} 0 ,
\end{multline*}
поскольку $f_k \to f$ в $\Lloc(\Om)$, а $\supp \varphi \Subset \Om$ — компакт.

\emph{Предел в правой части} получается аналогично:
\[
  \int\limits_{\supp \varphi} D^\alpha f_k \cdot \varphi\, dx
  \xrightarrow[k \to \infty]{}
  \int\limits_{\supp \varphi} g\, \varphi\, dx .
\]

Тогда перейдём к пределу:
\[
  \int\limits_{\supp \varphi} f\, D^\alpha \varphi\, dx
  = (-1)^{|\alpha|} \int\limits_{\supp \varphi} g\, \varphi\, dx
  \quad\Longrightarrow\quad
  \int\limits_\Om f\, D^\alpha \varphi\, dx
  = (-1)^{|\alpha|} \int\limits_\Om g\, \varphi\, dx ,
\]
то есть $g = D^\alpha f$ на $\Om$.
\end{proof}

\begin{idea}
Признак позволяет получать обобщённые производные предельным переходом: достаточно
подобрать удобную последовательность гладких (или просто более хороших) функций,
сходящуюся вместе с производными. Сходимость нужна лишь локально, в $\Lloc$, —
это самое слабое требование, какое ещё даёт результат.
\end{idea}

\begin{example}[$(|x|)' = \sgn x$]
Знаем, что $(|x|)' = \sgn x$ — проверим это признаком.
\end{example}

Возьмём $f = |x|$ и последовательность $f_k$, совпадающую с $|x|$ вне
$\left[-\tfrac1k, \tfrac1k\right]$ и сглаженную внутри этого отрезка.
Тогда $f_k \to f$, а производные $f_k'$ — кусочно-линейные функции, переходящие
от $-1$ к $1$ на $\left[-\tfrac1k, \tfrac1k\right]$, так что $f_k' \to \sgn x$.
По признаку~\ref{thm:crit-deriv} получаем то же самое: $(|x|)' = \sgn x$.

\inputfig{figures/fig-abs-x}

\subsection{Усреднение функций, имеющих обобщённые производные}

\begin{lemma}[Л1]\label{lem:avg-commute}
Пусть $u \in \Lp{1}(\Om)$ и $\exists\, D^\alpha u \in \Lp{1}(\Om)$. Тогда
\[
  (D^\alpha u)_\rho(x) = \big(D^\alpha u_\rho\big)(x) ,
\]
где $x \in \Om_\rho$, $\dist(\dOm, x) > 0$, а $\rho$ таково, что
$0 < \rho < \dist(\dOm, x)$.
\end{lemma}

\begin{idea}
Усреднение и обобщённое дифференцирование коммутируют — но только вдали от границы:
шар $\ball{x}{\rho}$ должен целиком лежать в $\Om$, иначе ядро «высунется» наружу,
где $u$ не определена.
\end{idea}

\begin{proof}
Проводим выкладку для $\partial/\partial x_i$; общий случай аналогичен.

При фиксированном $x$ функция $y \mapsto \omega_\rho(x-y)$ обращается в нуль вне
$\ball{x}{\rho}$, то есть финитна в $\Om$. Дифференцируя под знаком интеграла
(свойство 3 усреднения),
\[
  \dd{u_\rho}{x_i}(x)
  = \dd{}{x_i}\left( \int\limits_\Om \omega_\rho(x-y)\, u(y)\, dy \right)
  = \int\limits_\Om \dd{}{x_i}\,\omega_\rho(x-y)\, u(y)\, dy .
  \tag{$\ast\ast$}\label{eq:dstar}
\]
Далее, ядро зависит от разности, поэтому
\[
  \dd{}{x_i}\big(\omega_\rho(x-y)\big) = -\,\dd{}{y_i}\big(\omega_\rho(x-y)\big) ,
\]
и, подставляя это в \eqref{eq:dstar} и применяя определение обобщённой производной
к пробной функции $\varphi \equiv \omega_\rho(x - \cdot)$,
\[
  \eqref{eq:dstar} = -\int\limits_\Om \dd{\omega_\rho(x-y)}{y_i}\, u(y)\, dy
  = \int\limits_\Om \omega_\rho(x-y)\, \dd{u}{y_i}(y)\, dy
  = \left( \dd{u}{x_i} \right)_{\!\rho}(x) .
\]

Для произвольного мультииндекса всё то же самое, только знак набегает $|\alpha|$ раз:
\[
  \boxed{\;D^\alpha_x\, \omega_\rho(x-y) = (-1)^{|\alpha|}\, D^\alpha_y\, \omega_\rho(x-y)\;}
\]
\end{proof}

\begin{lemma}[Л2]\label{lem:avg-conv-loc}
Пусть $1 \le p < \infty$, $u \in \Lp{p}(\Om)$ и $\exists\, D^\alpha u \in \Lp{p}(\Om)$.
Тогда для любого $\Om' \Subset \Om$
\[
  u_\rho \to u \ \text{ в } \Lp{p}(\Om'), \qquad
  D^\alpha u_\rho \to D^\alpha u \ \text{ в } \Lp{p}(\Om') .
\]
\end{lemma}

\begin{warn}
В тетради во второй строке заключения сначала было написано $\Lp{p}(\Om)$ — зачёркнуто
и исправлено на $\Lp{p}(\Om')$. Сходимость именно локальная: около границы лемма~\ref{lem:avg-commute}
неприменима.
\end{warn}

\begin{proof}
Знаем, что $u_\rho \to u$ в $\Lp{p}(\Om)$ (теорема~\ref{thm:conv}), а значит и
$u_\rho \to u$ в $\Lp{p}(\Om')$.

Так же $(D^\alpha u)_\rho \to D^\alpha u$ в $\Lp{p}(\Om')$ — это та же теорема,
применённая к функции $D^\alpha u \in \Lp{p}(\Om)$.

Пусть $\dist(\Om', \dOm) = \rho_0 > 0$ и $0 < \rho < \rho_0$. Тогда по
лемме~\ref{lem:avg-commute} в $\Om'$ выполнено $(D^\alpha u)_\rho = D^\alpha u_\rho$,
откуда
\[
  D^\alpha u_\rho \to D^\alpha u \ \text{ в } \Lp{p}(\Om') . \qedhere
\]
\end{proof}

\begin{lemma}[Л3]\label{lem:const}
Пусть $u \in \Lloc(\Om)$, область $\Om$ не обязательно ограничена, и $Du = 0$
(все обобщённые производные первого порядка равны нулю). Тогда $u = \mathrm{const}$
почти всюду в $\Om$.
\end{lemma}

\begin{proof}
\textbf{1) Случай $u \in \Lp{1}(\Om)$, $\Om$ ограничена.}

Для $\eps > 0$ положим $\Om_\eps = \{x \in \Om \mid \dist(x, \dOm) > \eps\}$.
По лемме~\ref{lem:avg-commute} при $0 < \rho < \eps$ в $\Om_\eps$
\[
  \left( \dd{u}{x_i} \right)_{\!\rho} = \dd{u_\rho}{x_i} ,
  \qquad\text{причём}\qquad
  \left( \dd{u}{x_i} \right)_{\!\rho}
  = \int\limits_{\Om_\eps} \omega_\rho(x-y)\, \dd{u}{x_i}(y)\, dy = 0 ,
\]
так как $\partial u/\partial x_i = 0$. Значит, все частные производные гладкой функции
$u_\rho$ в $\Om_\eps$ нулевые, и
\[
  u_\rho(x) = C(\rho) \quad \text{в } \Om_\eps ,
\]
поскольку $u_\rho$ — гладкая функция.

\emph{Покажем, что $C(\rho) < \infty$.} По теореме~\ref{thm:norm}
\[
  |C(\rho)|\cdot \mu(\Om_\eps) = \int\limits_{\Om_\eps} |u_\rho|\, dx
  \le \int\limits_\Om |u_\rho|\, dx \le \int\limits_\Om |u|\, dx < \infty ,
\]
откуда
\[
  |C(\rho)| \le \frac{1}{\mu(\Om_\eps)} \int\limits_\Om |u|\, dx < \infty .
\]
Пусть $C(\rho_k) \to C_0$ при $\rho \to 0$ (выбрали сходящуюся подпоследовательность).
Тогда
\begin{multline*}
  \int\limits_{\Om_\eps} |u - C_0|\, dx
  = \int\limits_{\Om_\eps} |u - C(\rho) + C(\rho) - C_0|\, dx \\
  \le \int\limits_{\Om_\eps} |u - u_\rho|\, dx + \int\limits_{\Om_\eps} |C(\rho) - C_0|\, dx
  \le \underbrace{\int\limits_{\Om_\eps} |u - u_\rho|\, dx}_{\to\, 0}
    + \underbrace{|C(\rho) - C_0|\cdot \mu(\Om_\eps)}_{\to\, 0} \longrightarrow 0 ,
\end{multline*}
то есть $\int_{\Om_\eps} |u - C_0|\, dx = 0$, а значит $u = C_0$ п.\,в. в $\Om_\eps$.

Положим $\eps = \tfrac1k$, то есть $\Om_\eps = \{x \in \Om \mid \dist(x,\dOm) > \tfrac1k\}$.
Тогда $\mu\big(\Om_{1/k}[u \neq C_0]\big) = 0$ и
\[
  \Om[u \neq C_0] = \bigcup_{k \in \N} \Om_{1/k}[u \neq C_0]
  \quad\Longrightarrow\quad u = C_0 \ \text{п.\,в. на } \Om .
\]

\textbf{2) Случай $u \in \Lloc(\Om)$, $\Om$ не обязательно ограничена.}

Для любого $\Om' \Subset \Om$ имеем $u \in \Lp{1}(\Om')$, а для этого случая всё уже
доказано, так что $u = \mathrm{const}$ в $\Om'$. Возьмём
\[
  \Om_k = \Big\{x \in \Om \ \Big|\ \dist(x, \dOm) > \tfrac1k \Big\} \cap \ball{0}{k},
  \qquad \Om_k \Subset \Om, \qquad \Om = \bigcup_{k \in \N} \Om_k .
\]
Тогда $\mu\big(\Om_k[u \neq C_0]\big) = 0$, а
\[
  \Om[u \neq C_0] = \bigcup_{k=1}^{\infty} \Om_k[u \neq C_0]
  \quad\Longrightarrow\quad \mu\big(\Om[u \neq C_0]\big) = 0 . \qedhere
\]
\end{proof}

\begin{idea}
Пересечение с шаром $\ball{0}{k}$ во втором пункте нужно именно потому, что $\Om$ может
быть неограниченной: одно условие $\dist(x,\dOm) > 1/k$ компактности не даёт.
\end{idea}

\begin{theorem}[замена переменных]\label{thm:change-vars}
Пусть $\Om, \Om' \subset \R^n$, а $\varphi : \Om' \to \Om$ — дифференцируемая биекция
с непрерывными первыми производными и ненулевым якобианом. Пусть
\[
  u \in \Lloc(\Om), \qquad \exists\, \dd{u}{x_i} \in \Lloc(\Om), \quad i = \overline{1,n} .
\]
Тогда $v(y) = u(\varphi(y)) \in \Lloc(\Om')$ и $v$ имеет обобщённые производные
из $\Lloc(\Om')$:
\[
  \dd{v}{y_i} = \sum_{k=1}^{n} \dd{u}{x_k}\big(\varphi(y)\big)\, \dd{\varphi_k(y)}{y_i} .
\]
\end{theorem}

\begin{proof}
\plusplus
\end{proof}

\subsection{Случай одной независимой переменной}

\begin{definition}
Функция $u : (a,b) \to \R$ называется \textbf{абсолютно непрерывной}, если
\[
  \exists\, v \in \Lp{1}(a,b): \qquad u(x) = \int\limits_a^x v(t)\, dt + C .
\]
\end{definition}

\begin{remark}
$u$ абсолютно непрерывна $\Longrightarrow$ $u$ непрерывна.
\end{remark}

\begin{warn}
В тетради теорема ниже сначала была записана в виде «для того, чтобы $u \in \Lp{1}(a,b)$
имела первую обобщённую производную, необходимо и достаточно, чтобы $u$ была абсолютно
непрерывна», и её доказательство занимало две страницы — \textbf{обе перечёркнуты
крест-накрест}. Затем формулировка и доказательство переписаны заново; ниже приведён
именно второй, итоговый вариант. В зачёркнутой версии доказательство
необходимости шло через $W(x) = \int_a^x \frac{du}{dt}(t)\,dt$ и равенство
$\frac{d(W-u)}{dx} = 0 \Rightarrow W = u + C$ п.\,в.
\end{warn}

\begin{theorem}\label{thm:ac}
Функция $u \in \Lp{1}(a,b)$ абсолютно непрерывна $\iff$ существует
$u' \in \Lp{1}(a,b)$ (здесь $u'$ — соболевская производная).
\end{theorem}

\begin{proof}
\textbf{$\Longrightarrow$.} Пусть $u \in \Lp{1}(a,b)$ абсолютно непрерывна, то есть
\[
  \exists\, v \in \Lp{1}(a,b): \quad u(x) = \int\limits_a^x v(t)\, dt + C .
\]
Пользуемся признаком существования производной (теорема~\ref{thm:crit-deriv}).
Гладкие функции плотны в $\Lp{1}$, поэтому
\[
  \exists\, v_m \in C^\infty\big((a,b)\big): \quad v_m \to v \ \text{ в } \Lp{1}(a,b) .
\]
Положим
\[
  u_m(x) = \int\limits_a^x v_m(t)\, dt + C .
\]
Тогда для всех $x \in [a,b]$
\[
  |u_m(x) - u(x)| \le \int\limits_a^x |v_m(t) - v(t)|\, dt
  \le \int\limits_a^b |v_m(t) - v(t)|\, dt = \norm{v_m - v}_{\Lp{1}(a,b)} ,
\]
то есть $u_m \unif u$ на $(a,b)$, а значит $u_m \to u$ в $\Lp{1}(a,b)$.

Так как $u_m, v_m \in C^\infty$, то $u_m' = v_m$, причём
\[
  u_m \to u, \qquad v_m \to v \qquad \text{в } \Lp{1}(a,b) .
\]
Тогда по признаку существования производной $u' = v \in \Lp{1}(a,b)$.

\textbf{$\Longleftarrow$.} Пусть $u \in \Lp{1}(a,b)$ и существует $u' = w \in \Lp{1}(a,b)$.
Положим
\[
  \widetilde u(x) = \int\limits_a^x w(t)\, dt + C ,
\]
тогда $\widetilde u$ абсолютно непрерывна, и по уже доказанной части
$\widetilde u\,' = w$. Отсюда
\[
  (u - \widetilde u)' = w - w = 0
  \quad\overset{\text{Л\ref{lem:const}}}{\Longrightarrow}\quad
  u = \widetilde u + C_1
  \quad\Longrightarrow\quad
  u(x) = \int\limits_a^x w(t)\, dt + \underbrace{C + C_1}_{C_2} ,
\]
то есть $u$ абсолютно непрерывна.
\end{proof}

\subsubsection*{Почему это называется абсолютной непрерывностью}

Если $u \in C\big([a,b]\big)$, то $u$ равномерно непрерывна:
\[
  \forall \eps > 0 \ \exists \delta > 0: \quad |x - y| < \delta
  \ \Longrightarrow\ |u(x) - u(y)| < \eps .
\]
Классическое определение абсолютной непрерывности требует большего — контроля не над
одним отрезком, а над любым конечным набором непересекающихся отрезков:
\begin{idea}[title={Классическое определение}]
Для $a \le x_1 < y_1 < x_2 < y_2 < \dots < x_k \le y_k \le b$
\[
  \sum_{i=1}^{k} (x_i - y_i) < \delta
  \quad\Longrightarrow\quad
  \sum_{i=1}^{k} |u(x_i) - u(y_i)| < \eps .
\]
\end{idea}

\begin{warn}
В тетради в сумме написано $\sum (x_i - y_i)$ при $x_i < y_i$, то есть формально
отрицательные слагаемые. Имеется в виду суммарная длина отрезков
$\sum (y_i - x_i) < \delta$.
\end{warn}

\subsubsection*{Канторова лестница}

\inputfig{figures/fig-cantor}

Канторова лестница на самом канторовом множестве не задана явно, но продолжается
на весь отрезок по непрерывности. Она непрерывна, однако суммарная длина отрезков
$m$-го уровня равна $(2/3)^m \to 0$, тогда как суммарное приращение функции на них
остаётся равным $1$. Значит, \textbf{канторова лестница не абсолютно непрерывна}, хотя
«если абсолютно непрерывна, то непрерывна».

\begin{idea}
Это и показывает, что абсолютная непрерывность строго сильнее равномерной: именно она,
а не непрерывность, равносильна существованию суммарной производной.
\end{idea}

\begin{proposition}[эквивалентность определений]
Если $u$ абсолютно непрерывна в смысле определения через интеграл, то она абсолютно
непрерывна и в классическом смысле, и наоборот.
\end{proposition}

\begin{proof}
\plusplus
\end{proof}

\lecture{Лекция 5}{06.10.2026}

\subsection{Определения и общие свойства пространств Соболева}

\begin{definition}
\textbf{Пространство Соболева} $\Wp{p}{\ell}(\Om)$, где $\ell \in \N$, $1 \le p \le \infty$,
составляют функции, которые имеют обобщённые производные до порядка $\ell$ включительно
и суммируемые вместе с этими производными со степенью $p$, то есть
\[
  u \in \Wp{p}{\ell}(\Om) \iff \exists\, D^\alpha u \in \Lp{p}(\Om)
  \quad \forall\, |\alpha| \le \ell .
\]
\end{definition}

\begin{remark}
$\Wp{p}{0}(\Om) = \Lp{p}(\Om)$ по определению.
\end{remark}

\begin{definition}[норма в $\Wp{p}{\ell}$]
\[
  \norm{u}^{(1)}_{\Wp{p}{\ell}(\Om)}
  = \sum_{|\alpha| \le \ell} \norm{D^\alpha u}_{\Lp{p}(\Om)} .
\]
Если $1 \le p < \infty$, то можно ввести ещё одну норму:
\[
  \norm{u}^{(2)}_{\Wp{p}{\ell}(\Om)}
  = \left( \sum_{|\alpha| \le \ell} \int\limits_\Om |D^\alpha u|^p\, dx \right)^{1/p} .
\]
\end{definition}

\begin{warn}
В тетради обе нормы обозначены одинаково, $\norm{u}_{\Wp{p}{\ell}(\Om)}$, и дальше
в доказательстве различаются припиской $\norm{u}^1$ и $\norm{u}^2$. Здесь они с самого
начала помечены верхними индексами $(1)$ и $(2)$. Осторожно: в тетради
\emph{сумма норм} обозначена как $\norm{u}^2$, а \emph{корень из суммы $p$-х степеней} —
как $\norm{u}^1$; ниже сохранено именно это соответствие.
\end{warn}

\begin{proposition}\label{prop:norm-equiv}
При $1 \le p \le \infty$ нормы эквивалентны: $\norm{u}^{(1)} \sim \norm{u}^{(2)}$.
\end{proposition}

\begin{proof}
Надо указать $c_1, c_2$ такие, что
\[
  c_1 \norm{u}^{(1)} \le \norm{u}^{(2)} \le c_2 \norm{u}^{(1)}
  \qquad \forall u \in \Wp{p}{\ell}(\Om) .
\]

\begin{reminder}[неравенства, которыми пользуемся]
Минковского: $\displaystyle \left(\int |f+g|^p\right)^{1/p}
\le \left(\int |f|^p\right)^{1/p} + \left(\int |g|^p\right)^{1/p}$.

\smallskip
Гёльдера: $\displaystyle \int fg\, dx
\le \left(\int |f|^p dx\right)^{1/p} \left(\int |g|^q dx\right)^{1/q}$, в частности
\[
  \int\limits_\Om |D^\alpha u|\, dx
  \le \left( \int\limits_\Om |D^\alpha u|^p\, dx \right)^{1/p}
      \underbrace{\left( \int\limits_\Om 1^q\, dx \right)^{1/q}}_{=\,(\mu(\Om))^{1/q}},
  \qquad \frac1p + \frac1q = 1 .
\]
\end{reminder}

Обозначим
\[
  a_\alpha = \norm{D^\alpha u}_{\Lp{p}(\Om)} \ge 0 ,
\]
а через $N$ — число мультииндексов $|\alpha| \le \ell$. Тогда
\[
  \norm{u}^{(1)} = \left( \sum_{|\alpha| \le \ell} a_\alpha^p \right)^{1/p},
  \qquad
  \norm{u}^{(2)} = \sum_{|\alpha| \le \ell} a_\alpha ,
  \qquad c_1 = 1, \quad c_2 = N^{1 - \frac1p} = N^{\frac1q} .
\]

\textbf{1) $c_1 \norm{u}^{(1)} \le \norm{u}^{(2)}$.} Положим
$S = \sum_{|\alpha| \le \ell} a_\alpha = \norm{u}^{(2)}$. Тогда $a_\beta \le S$ для
всех $\beta$ с $|\beta| \le \ell$, откуда
\[
  a_\beta^p = a_\beta \cdot a_\beta^{\,p-1} \le a_\beta\, S^{p-1}
  \quad\Longrightarrow\quad
  \sum_{|\beta| \le \ell} a_\beta^p \le S^{p-1} \underbrace{\sum_{|\beta| \le \ell} a_\beta}_{=\,S} = S^p ,
\]
то есть $\big(\norm{u}^{(1)}\big)^p \le \big(\norm{u}^{(2)}\big)^p$ и
$\norm{u}^{(1)} \le \norm{u}^{(2)}$.

\textbf{2) $\norm{u}^{(2)} \le N^{1-\frac1p} \norm{u}^{(1)}$.}
При $p = 1$ показатель $N^{1-1/p} = N^0 = 1$, и неравенство тривиально. При $p > 1$
берём сопряжённый показатель $q = \frac{p}{p-1} = \left(1 - \frac1p\right)^{-1}$
и применяем Гёльдера к сумме:
\[
  \norm{u}^{(2)} = \sum_{|\alpha| \le \ell} a_\alpha \cdot 1
  \le \left( \sum_{|\alpha| \le \ell} a_\alpha^p \right)^{1/p}
      \left( \sum_{|\alpha| \le \ell} 1^q \right)^{1/q}
  = \norm{u}^{(1)} \cdot N^{1/q} ,
\]
то есть $\norm{u}^{(2)} \le c_2 \norm{u}^{(1)}$.
\end{proof}

\begin{warn}
На доске в пункте~2 написано «$p = 1 \Rightarrow q = 0$». Сопряжённый показатель при
$p=1$ равен $\infty$, а не нулю; существенно здесь другое — что множитель
$N^{1-1/p}$ при $p=1$ обращается в единицу.
\end{warn}

\begin{theorem}[Т1]\label{thm:w-banach}
$\Wp{p}{\ell}(\Om)$ — банахово пространство.
\end{theorem}

\begin{proof}
Докажем только для $1 \le p < \infty$.

Пусть $\{u_n\} \subset \Wp{p}{\ell}(\Om)$ — фундаментальная последовательность:
\[
  \sum_{|\alpha| \le \ell} \norm{D^\alpha u_n - D^\alpha u_m}_{\Lp{p}(\Om)}
  \xrightarrow[n,m \to \infty]{} 0
\]
\[
  \iff \norm{D^\alpha u_n - D^\alpha u_m}_{\Lp{p}(\Om)} \xrightarrow[n,m \to \infty]{} 0
  \quad \forall \alpha: |\alpha| \le \ell
\]
\[
  \iff \{D^\alpha u_n\} \ \text{— фундаментальна в } \Lp{p}(\Om)
  \quad \forall \alpha: |\alpha| \le \ell
\]
\[
  \iff \forall \alpha: |\alpha| \le \ell \quad \exists\, u_\alpha \in \Lp{p}(\Om):
  \quad D^\alpha u_n \to u_\alpha \ \text{ в } \Lp{p}(\Om) ,
\]
последнее — в силу полноты $\Lp{p}(\Om)$.

Обозначим $u = u_0$ (предел самих $u_n$). Имеем
\[
  u_n \to u, \qquad D^\alpha u_n \to u_\alpha \quad \forall \alpha ,
\]
и по признаку существования обобщённой производной (теорема~\ref{thm:crit-deriv})
\[
  \exists\, D^\alpha u \in \Lp{p}(\Om): \quad D^\alpha u = u_\alpha .
\]

\begin{idea}
Признак сформулирован для сходимости в $\Lloc$, а здесь сходимость идёт в $\Lp{p}(\Om)$.
Это не мешает: на любом $\Om' \Subset \Om$ из сходимости в $\Lp{p}$ следует сходимость
в $\Lp{1}(\Om')$ (неравенство Гёльдера на множестве конечной меры), то есть сходимость
в $\Lloc(\Om)$. Именно это замечание и стоит на полях в тетради.
\end{idea}

Тогда $D^\alpha u_n \to D^\alpha u$ в $\Lp{p}(\Om)$ для всех $|\alpha| \le \ell$, что
по определению нормы означает $u_n \to u$ в $\Wp{p}{\ell}(\Om)$.
\end{proof}

\subsubsection*{Пространства $\big[\Lp{p}(\Om)\big]^N$}

\begin{idea}
Цель конструкции — представить $\Wp{p}{\ell}(\Om)$ как подпространство прямой степени
$\Lp{p}$, сопоставив функции набор всех её производных. Сепарабельность и рефлексивность
тогда наследуются «сверху».
\end{idea}

Пусть $X$ — сепарабельное банахово пространство и
\[
  X^N = \{(x_1, x_2, \dots, x_N)\}, \qquad x_i \in X, \quad i = \overline{1,N} ,
\]
причём для вектора $\underline{x} = (x_1, \dots, x_N)$
\[
  \norm{\underline{x}}_{X^N} = \left( \sum_{i=1}^{N} \norm{x_i}_X^p \right)^{1/p},
  \qquad 1 \le p < \infty .
\]

\begin{proposition}\label{prop:xn}
Тогда $X^N$ — сепарабельное банахово пространство.
\end{proposition}

\begin{proof}
\emph{Сепарабельность.} Пусть $\underline{x} = (x_1, \dots, x_N)$. Так как $X$
сепарабельно, существует счётное всюду плотное $A \subset X$, и для любого $\eps > 0$
найдутся номера $n_j$ с $a_{n_j} \in A$ такие, что
\[
  \norm{x_j - a_{n_j}}_X < \frac{\eps}{N^{1/p}} .
\]
Пусть $\underline{a} = (a_{n_1}, \dots, a_{n_N}) \in A^N$. Тогда
\[
  \norm{\underline{x} - \underline{a}}_{X^N}
  = \Bigg( \sum_{j=1}^{N} \underbrace{\norm{x_j - a_{n_j}}_X^p}_{<\ \eps^p/N} \Bigg)^{1/p} < \eps .
\]
То есть для любых $\underline{x} \in X^N$ и $\eps > 0$ найдётся $\underline{a} \in A^N$
с $\norm{\underline{x} - \underline{a}}_{X^N} < \eps$, а $A^N$ счётно.

\emph{Полнота.} Пусть $\{\underline{x}^i\}$ — фундаментальная последовательность в $X^N$:
\[
  \norm{\underline{x}^i - \underline{x}^j}_{X^N}
  = \left( \sum_{k=1}^{N} \norm{x_k^i - x_k^j}_X^p \right)^{1/p}
  \xrightarrow[i,j \to \infty]{} 0
  \iff \norm{x_k^i - x_k^j}_X \xrightarrow[i,j \to \infty]{} 0 \quad \forall k = \overline{1,N}
\]
\[
  \iff \{x_k^i\}_{i=1}^{\infty} \ \text{— фундаментальна в } X \ \ \forall k
  \iff x_k^i \xrightarrow[i \to \infty]{} x_k \in X \quad \forall k = \overline{1,N} .
\]
Положив $\underline{x}_0 = (x_1, \dots, x_N) \in X^N$, получаем
$\underline{x}^i \to \underline{x}_0$ при $i \to \infty$.
\end{proof}

Применяем это к $X = \Lp{p}(\Om)$:
\[
  \big[\Lp{p}(\Om)\big]^N
  = \big\{\underline{u} = (u_1, \dots, u_N) \ \big|\ u_i \in \Lp{p}(\Om) \big\} ,
\]
\[
  \norm{\underline{u}}_{[\Lp{p}(\Om)]^N}
  = \left( \sum_{i=1}^{N} \norm{u_i}^p_{\Lp{p}(\Om)} \right)^{1/p}
  = \left( \sum_{i=1}^{N} \int\limits_\Om |u_i|^p\, dx \right)^{1/p},
  \qquad 1 \le p < \infty .
\]
Пространство $\Lp{p}(\Om)$ полное и сепарабельное, поэтому $\big[\Lp{p}(\Om)\big]^N$
тоже полное и сепарабельное.

\begin{proposition}[сопряжённое пространство]\label{prop:dual}
Пусть $1 < p < \infty$ и $p' = \frac{p}{p-1}$.
\begin{enumerate}
  \item Для любого $\underline{v} \in \big[\Lp{p'}(\Om)\big]^N$ формула
        \[
          \ell(\underline{u}) = \sum_{i=1}^{N} \int\limits_\Om u_i v_i\, dx ,
          \qquad \underline{u} \in \big[\Lp{p}(\Om)\big]^N ,
        \]
        задаёт линейный ограниченный функционал
        $\ell \in \big(\big[\Lp{p}(\Om)\big]^N\big)^*$, причём
        \[
          \norm{\ell}_{([\Lp{p}(\Om)]^N)^*} = \norm{\underline{v}}_{[\Lp{p'}(\Om)]^N} .
        \]
  \item Обратно, для любого функционала
        $\ell \in \big(\big[\Lp{p}(\Om)\big]^N\big)^*$ существует единственный
        $\underline{v} \in \big[\Lp{p'}(\Om)\big]^N$, для которого выполнены оба
        равенства пункта~1.
\end{enumerate}
То есть существует изометрический изоморфизм
\[
  \big(\big[\Lp{p}(\Om)\big]^N\big)^* \simeq \big[\Lp{p'}(\Om)\big]^N .
\]
\end{proposition}

\begin{proof}
\plusplus
\end{proof}

\begin{corollary}
Так как $(p')' = p$, то
\[
  \Big(\big(\big[\Lp{p}(\Om)\big]^N\big)^*\Big)^* \simeq \big[\Lp{p}(\Om)\big]^N ,
\]
то есть $\big[\Lp{p}(\Om)\big]^N$ рефлексивно.
\end{corollary}

\begin{theorem}[Т2]\label{thm:w-sep-refl}
$\Wp{p}{\ell}(\Om)$ сепарабельно при $1 \le p < \infty$ и рефлексивно при $1 < p < \infty$.
\end{theorem}

\begin{proof}
Пусть $N = \sum_{|\alpha| \le \ell} 1$ — число мультииндексов. Рассмотрим
пространство $\big[\Lp{p}(\Om)\big]^N$ и положим
\[
  \overline{W} = \Big\{ \underline{u} \in \big[\Lp{p}(\Om)\big]^N
  \ \Big|\ \underline{u} = (u_\alpha)_{|\alpha| \le \ell}:
  \ \exists\, u \in \Wp{p}{\ell}(\Om), \ u_\alpha = D^\alpha u \Big\} .
\]
Нужно связать $\big[\Lp{p}(\Om)\big]^N$ и $\overline{W}$.

\emph{Изоморфизм.} Пусть $u, v \in \Wp{p}{\ell}(\Om)$ и
$u \leftrightarrow \underline{u} \in \overline{W}$,
$v \leftrightarrow \underline{v} \in \overline{W}$. Тогда
\[
  a u + b v \in \Wp{p}{\ell}(\Om)
  \qquad\longleftrightarrow\qquad
  a \underline{u} + b \underline{v} \in \overline{W} ,
\]
то есть соответствие $\Wp{p}{\ell}(\Om) \leftrightarrow \overline{W}$ — биекция,
сохраняющая линейные операции, а это и есть изоморфизм.

\emph{Изометричность.}
\[
  \norm{\underline{u}}_{[\Lp{p}(\Om)]^N}
  = \left( \sum_{|\alpha| \le \ell} \norm{u_\alpha}^p_{\Lp{p}(\Om)} \right)^{1/p}
  = \left( \sum_{|\alpha| \le \ell} \norm{D^\alpha u}^p_{\Lp{p}(\Om)} \right)^{1/p}
  = \norm{u}_{\Wp{p}{\ell}(\Om)} ,
\]
то есть изоморфизм изометрический.

\emph{Подпространство.} $\Wp{p}{\ell}(\Om)$ банахово (теорема~\ref{thm:w-banach}),
значит и $\overline{W}$ — банахово пространство. При этом
$\overline{W} \subset \big[\Lp{p}(\Om)\big]^N$ — линейное многообразие, а линейное
многообразие является подпространством тогда и только тогда, когда оно банахово
относительно индуцированной нормы. Значит, $\overline{W}$ — подпространство
в $\big[\Lp{p}(\Om)\big]^N$.

Наконец, $\big[\Lp{p}(\Om)\big]^N$ сепарабельно и рефлексивно, а эти свойства
наследуются подпространствами, поэтому $\overline{W}$ сепарабельно и рефлексивно.
Вместе с $\overline{W} \simeq \Wp{p}{\ell}(\Om)$ это и даёт утверждение теоремы.
\end{proof}

\begin{theorem}[Т3, замена переменных в $\Wp{p}{1}$]\label{thm:w-change}
Пусть $\varphi : \omega \to \Om$ — диффеоморфизм (все производные самого отображения
и обратного ограничены), $|J(\varphi)| > 0$ на $\overline{\omega}$ и
$\left| \dd{\varphi_k}{y_i} \right| \le \varkappa$ для некоторого числа $\varkappa$.
Пусть $u \in \Wp{p}{1}(\Om)$ и $v(y) = u(\varphi(y))$. Тогда $v \in \Wp{p}{1}(\omega)$ и
\[
  \exists\, c_1, c_2 > 0: \qquad
  c_1 \norm{u}_{\Wp{p}{1}(\Om)} \le \norm{v}_{\Wp{p}{1}(\omega)}
  \le c_2 \norm{u}_{\Wp{p}{1}(\Om)} .
\]
\end{theorem}

\begin{remark}
Этот результат верен и для $\ell \in \N$, если потребовать побольше гладкости:
$\varphi \in C^\ell$.
\end{remark}

\begin{proof}
\plusplus
\end{proof}

\subsection{Пространства $\Wpz{p}{\ell}(\Om)$}

\begin{definition}
\[
  \Wpz{p}{\ell}(\Om) = \mathrm{Cl}_{\Wp{p}{\ell}(\Om)}\big(C_0^\infty(\Om)\big)
\]
— замыкание финитных функций по норме $\Wp{p}{\ell}(\Om)$. То есть
\[
  u \in \Wpz{p}{\ell}(\Om) \iff u \in \Wp{p}{\ell}(\Om):
  \ \exists\, \{u_n\} \subset C_0^\infty(\Om): \ u_n \xrightarrow[n \to \infty]{} u
  \ \text{ в } \Wp{p}{\ell}(\Om) .
\]
\end{definition}

\begin{remark}
$\Wpz{p}{\ell}(\Om)$ — подпространство $\Wp{p}{\ell}(\Om)$, поэтому оно сепарабельно,
рефлексивно и банахово при $1 < p < \infty$.
\end{remark}

\begin{proposition}[утв.\ 1: продолжение нулём]\label{prop:w0-ext}
Пусть $u \in \Wpz{p}{\ell}(\Om)$, $\Om \subset \widetilde\Om$ и
\[
  \widetilde u = \begin{cases} u, & \text{в } \Om, \\ 0, & \text{в } \widetilde\Om \setminus \Om. \end{cases}
\]
Тогда $\widetilde u \in \Wpz{p}{\ell}(\widetilde\Om)$ и
$\norm{u}_{\Wpz{p}{\ell}(\Om)} = \norm{\widetilde u}_{\Wpz{p}{\ell}(\widetilde\Om)}$.
\end{proposition}

\begin{proof}
Из $u \in \Wpz{p}{\ell}(\Om)$ следует существование $\{u_n\} \subset C_0^\infty(\Om)$
с $u_n \to u$ в $\Wp{p}{\ell}(\Om)$. Положим
\[
  \widetilde u_n = \begin{cases} u_n, & \text{в } \Om, \\ 0, & \text{в } \widetilde\Om \setminus \Om, \end{cases}
  \qquad\Longrightarrow\qquad \widetilde u_n \in C_0^\infty(\widetilde\Om) ,
\]
что законно: носитель $u_n$ — компакт строго внутри $\Om$, поэтому продолжение нулём
остаётся гладким. Так как
\[
  \norm{\widetilde u_n - \widetilde u_m}_{\Wp{p}{\ell}(\widetilde\Om)}
  = \norm{u_n - u_m}_{\Wp{p}{\ell}(\Om)} ,
\]
то из фундаментальности $\{u_n\}$ в $\Wp{p}{\ell}(\Om)$ следует фундаментальность
$\{\widetilde u_n\}$ в $\Wp{p}{\ell}(\widetilde\Om)$, а значит и в
$\Wpz{p}{\ell}(\widetilde\Om)$. Последнее — банахово, поэтому
\[
  \exists\, \widetilde u \in \Wpz{p}{\ell}(\widetilde\Om):
  \quad \widetilde u_n \to \widetilde u \ \text{ в } \Wpz{p}{\ell}(\widetilde\Om) .
\]

\emph{Вопрос: действительно ли $\widetilde u$ — это продолжение $u$ нулём?} Имеем
\begin{align*}
  \norm{\widetilde u_n - \widetilde u}_{\Wp{p}{\ell}(\widetilde\Om)}
  &= \sum_{|\alpha| \le \ell}
     \left( \int\limits_{\widetilde\Om} |D^\alpha(\widetilde u_n - \widetilde u)|^p\, dx \right)^{1/p}
     \xrightarrow[n \to \infty]{} 0 , \\
  \norm{\widetilde u_n - \widetilde u}_{\Wp{p}{\ell}(\widetilde\Om)}
  &= \sum_{|\alpha| \le \ell}
     \left( \int\limits_{\Om} |D^\alpha(u_n - \widetilde u)|^p\, dx
          + \int\limits_{\widetilde\Om \setminus \Om} |D^\alpha \widetilde u|^p\, dx \right)^{1/p} .
\end{align*}
Отбрасывая второе слагаемое, получаем
\[
  \norm{\widetilde u_n - \widetilde u}_{\Wp{p}{\ell}(\widetilde\Om)}
  \ge \sum_{|\alpha| \le \ell}
      \left( \int\limits_{\Om} |D^\alpha(u_n - \widetilde u)|^p\, dx \right)^{1/p}
  = \norm{u_n - \widetilde u}_{\Wp{p}{\ell}(\Om)} \longrightarrow 0 ,
\]
то есть $u_n \to \widetilde u$ в $\Wp{p}{\ell}(\Om)$. Но и $u_n \to u$ в
$\Wp{p}{\ell}(\Om)$, значит $u = \widetilde u$ в $\Om$. Аналогично (отбрасывая, наоборот,
первое слагаемое) получаем $\widetilde u = 0$ в $\widetilde\Om \setminus \Om$.
\end{proof}

\begin{proposition}[утв.\ 2: усреднения сходятся в норме $\Wp{p}{\ell}$]\label{prop:w0-avg}
Пусть $u \in \Wpz{p}{\ell}(\Om)$. Тогда $u_\rho \to u$ в $\Wp{p}{\ell}(\Om)$ при $\rho \to 0$.
\end{proposition}

\begin{proof}
Возьмём $\Om \Subset \widetilde\Om$ и $0 < \rho < \dist(\Om, \partial\widetilde\Om)$,
положив $u = 0$ в $\widetilde\Om \setminus \Om$ (утверждение~\ref{prop:w0-ext}). Тогда
\[
  u_\rho(x) = \int\limits_{\widetilde\Om} \omega_\rho(x-y)\, u(y)\, dy
            = \int\limits_{\Om} \omega_\rho(x-y)\, u(y)\, dy .
\]
Знаем, что $u_\rho \to u$ в $\Lp{p}(\widetilde\Om)$ (теорема~\ref{thm:conv}); здесь
существенно, что $u = 0$ в $\widetilde\Om\setminus\Om$, то есть $u \in \Lp{p}(\widetilde\Om)$.

Из леммы про усреднения (Л\ref{lem:avg-conv-loc})
\[
  D^\alpha u_\rho \xrightarrow[\rho \to 0]{} D^\alpha u \ \text{ в } \Lp{p}(\Om')
  \qquad \forall\, \Om' \Subset \widetilde\Om ,
\]
в частности для $\Om' = \Om$. Итого
\[
  u_\rho \to u \ \text{ в } \Lp{p}(\Om), \qquad
  D^\alpha u_\rho \to D^\alpha u \ \text{ в } \Lp{p}(\Om)
  \quad\Longrightarrow\quad
  u_\rho \to u \ \text{ в } \Wp{p}{\ell}(\Om) . \qedhere
\]
\end{proof}

\begin{remark}
Существенно, что $u \in \Wpz{p}{\ell}(\Om)$: для произвольной $u \in \Wp{p}{\ell}(\Om)$
так получалось бы не всегда.
\end{remark}

\begin{idea}
Причина — в утверждении~\ref{prop:w0-ext}: только для функций из $\Wpz{p}{\ell}$
продолжение нулём не портит обобщённые производные, а значит усреднение можно брать
по увеличенной области. У произвольной $u \in \Wp{p}{\ell}(\Om)$ продолжение нулём даёт
скачок на границе, и у $D^\alpha u_\rho$ около $\dOm$ сходимости нет.
\end{idea}

\begin{proposition}[утв.\ 3: финитность достаточна]\label{prop:w0-supp}
Пусть $u \in \Wp{p}{\ell}(\Om)$ и $\supp u \Subset \Om$. Тогда $u \in \Wpz{p}{\ell}(\Om)$.
\end{proposition}

\inputfig{figures/fig-supp-nested}

\begin{proof}
Надо предъявить $\{u_n\} \subset C_0^\infty(\Om)$ с $u_n \to u$ в $\Wp{p}{\ell}(\Om)$;
возьмём в качестве них усреднения.

Выберем $\Om'$ с $\supp u \Subset \Om' \Subset \Om$ и
$0 < \rho < \dist(\supp u, \partial\Om')$. Тогда
\[
  u_\rho(x) = \int\limits_\Om \omega_\rho(x-y)\, u(y)\, dy
            = \int\limits_{\supp u} \omega_\rho(x-y)\, u(y)\, dy \in C_0^\infty(\R^n) ,
\]
причём $u_\rho \equiv 0$ в $\Om \setminus \Om'$, то есть $u_\rho \in C_0^\infty(\Om)$.

По лемме~\ref{lem:avg-conv-loc} $D^\alpha u_\rho \to D^\alpha u$ в $\Lp{p}(\Om')$ для всех
$|\alpha| \le \ell$, а вне $\Om'$ обе части равны нулю: $D^\alpha u_\rho = D^\alpha u = 0$.
Значит, $D^\alpha u_\rho \to D^\alpha u$ в $\Lp{p}(\Om)$, то есть $u_\rho \to u$
в $\Wp{p}{\ell}(\Om)$, и по определению $u \in \Wpz{p}{\ell}(\Om)$.
\end{proof}

\subsubsection*{Пример: $\Wpz{p}{1}$ строго уже, чем $\Wp{p}{1}$}

\begin{example}
Пусть $\Om = (0,1)$ и $u \equiv 1$ на $\Om$. Тогда $u \in \Wp{p}{1}(\Om)$. Докажем,
что $u \notin \Wpz{p}{1}(\Om)$.
\end{example}

\inputfig{figures/fig-unit-ext}

\begin{proof}
Пусть, от противного, $u \in \Wpz{p}{1}(\Om)$. Возьмём $\widetilde\Om = (-1,2)$. По
утверждению~\ref{prop:w0-ext} продолжение нулём
\[
  \widetilde u = \begin{cases}
    0, & x \in (-1, 0], \\
    u, & x \in (0, 1), \\
    0, & x \in [1, 2),
  \end{cases}
\]
лежит в $\Wpz{p}{1}(\widetilde\Om) \subset \Wp{p}{1}(\widetilde\Om)$. Но тогда по
теореме~\ref{thm:ac} функция $\widetilde u$ абсолютно непрерывна, а значит непрерывна —
чего быть не может: у неё скачки в точках $0$ и $1$. Противоречие, следовательно
$u \notin \Wpz{p}{1}(\Om)$.
\end{proof}

\subsubsection*{Интегрирование по частям}

\begin{theorem}[интегрирование по частям]\label{thm:parts}
Пусть $u \in \Wp{p}{\ell}(\Om)$, $v \in \Wpz{p^*}{\ell}(\Om)$, где
$p^* = \frac{p}{p-1}$ и $1 < p < \infty$. Тогда
\[
  \int\limits_\Om D^\alpha u \cdot v\, dx
  = (-1)^{|\alpha|} \int\limits_\Om u\, D^\alpha v\, dx ,
  \qquad |\alpha| \le \ell .
\]
\end{theorem}

\begin{proof}
Возьмём $\{v_n\} \subset C_0^\infty(\Om)$ с $v_n \to v$ в $\Wp{p^*}{\ell}(\Om)$. Для
каждого $v_n$ равенство верно по определению обобщённой производной (ведь
$v_n \in C_0^\infty(\Om)$):
\[
  \int\limits_\Om D^\alpha u \cdot v_n\, dx
  = (-1)^{|\alpha|} \int\limits_\Om u\, D^\alpha v_n\, dx .
\]
Перейдём в каждой части к пределу.

\emph{Левая часть.} По неравенству Гёльдера
\begin{multline*}
  \left| \int\limits_\Om D^\alpha u \cdot v_n\, dx - \int\limits_\Om D^\alpha u \cdot v\, dx \right|
  \le \int\limits_\Om |D^\alpha u| \cdot |v_n - v|\, dx \\
  \le \left( \int\limits_\Om |D^\alpha u|^p\, dx \right)^{1/p}
      \Bigg( \underbrace{\int\limits_\Om |v_n - v|^{p^*}\, dx}_{\to\, 0} \Bigg)^{1/p^*}
  \xrightarrow[n \to \infty]{} 0 .
\end{multline*}

\emph{Правая часть.} Так же,
\begin{multline*}
  \left| \int\limits_\Om u\, D^\alpha v_n\, dx - \int\limits_\Om u\, D^\alpha v\, dx \right|
  \le \int\limits_\Om |u| \cdot |D^\alpha v_n - D^\alpha v|\, dx \\
  \le \left( \int\limits_\Om |u|^p\, dx \right)^{1/p}
      \Bigg( \underbrace{\int\limits_\Om |D^\alpha v_n - D^\alpha v|^{p^*}\, dx}_{\to\, 0} \Bigg)^{1/p^*}
  \longrightarrow 0 ,
\end{multline*}
поскольку $v_n \to v$ в $\Wp{p^*}{\ell}(\Om)$.
\end{proof}

\begin{idea}
Финитность $v$ нужна ровно для того, чтобы не возникло граничных членов: у гладких
финитных $v_n$ их нет по определению обобщённой производной, а предельный переход
в $\Wp{p^*}{\ell}$ их не возвращает. Для произвольной $v \in \Wp{p^*}{\ell}(\Om)$ формула,
вообще говоря, неверна — появляются следы на $\dOm$.
\end{idea}


\end{document}
